% =====================================================================
% Kapitel 5 - Digitale Signalverarbeitung und Orchestrierung (Python)
% Autor: Cem Acar
%
% EINBINDEN ins Hauptdokument:
%   1. Dieses File nach latexDocs/ legen.
%   2. Im Master (Dokumentation.tex) an der richtigen Stelle einfuegen:
%        % =====================================================================
% Kapitel 5 - Digitale Signalverarbeitung und Orchestrierung (Python)
% Autor: Cem Acar
%
% EINBINDEN ins Hauptdokument:
%   1. Dieses File nach latexDocs/ legen.
%   2. Im Master (Dokumentation.tex) an der richtigen Stelle einfuegen:
%        % =====================================================================
% Kapitel 5 - Digitale Signalverarbeitung und Orchestrierung (Python)
% Autor: Cem Acar
%
% EINBINDEN ins Hauptdokument:
%   1. Dieses File nach latexDocs/ legen.
%   2. Im Master (Dokumentation.tex) an der richtigen Stelle einfuegen:
%        % =====================================================================
% Kapitel 5 - Digitale Signalverarbeitung und Orchestrierung (Python)
% Autor: Cem Acar
%
% EINBINDEN ins Hauptdokument:
%   1. Dieses File nach latexDocs/ legen.
%   2. Im Master (Dokumentation.tex) an der richtigen Stelle einfuegen:
%        \input{Kapitel5_DSP_Cem}
%   3. Der Master braucht nur diese Pakete (Dokumentation.tex hat beide schon):
%        \usepackage{amsmath}
%        \usepackage{listings}
%      Kein xcolor/babel noetig: Umlaute laufen unter UTF-8, deutsche
%      Anfuehrungszeichen sind literal „..." gesetzt.
%
% STANDALONE-VORSCHAU (zum alleine Kompilieren einfach drumherum setzen):
%   \documentclass{article}\usepackage{amsmath}\usepackage{listings}
%   \begin{document}
%   \input{Kapitel5_DSP_Cem}
%   \end{document}
% =====================================================================

% Listing-Stil fuer dieses Kapitel (nur listings-Kern, keine Farben noetig)
\lstdefinestyle{cemdsp}{
    language=Python,
    basicstyle=\ttfamily\small,
    showstringspaces=false,
    breaklines=true,
    numbers=left,
    numberstyle=\tiny,
    frame=single,
    tabsize=4
}

%======================================================================
\section{Digitale Signalverarbeitung und Orchestrierung (Python)}
\label{sec:python-dsp}

Die Rust-Engine berechnet die Physik der Schallausbreitung, der Python-Dienst
übernimmt die Signalverarbeitung und die Orchestrierung der Pipeline. Warum die
Physik in Rust und die Signalverarbeitung in Python mit NumPy und SciPy liegt,
ist in Kapitel~3 begründet. Die signaltheoretischen Grundlagen, also
Impulsantwort, Faltungsintegral und schnelle Faltung über die FFT, behandelt
Abschnitt~2.3. Dieses Kapitel zeigt, wie diese Grundlagen im Code umgesetzt sind,
und folgt dabei dem Datenfluss durch die Pipeline.

\subsection{Pipeline-Management durch \texttt{main.py}}
\label{subsec:pipeline}

\texttt{main.py} führt die vier Verarbeitungsschritte nacheinander aus. Tritt in
einem Schritt ein Fehler auf, bricht die Pipeline sofort mit einer Meldung ab,
statt mit ungültigen Zwischendaten weiterzurechnen (fail-fast).

\begin{enumerate}
    \item \textbf{Physik-Engine starten:} \texttt{run\_rust\_and\_load\_json} ruft
    die vorkompilierte Rust-Binary als Subprozess auf und lädt deren Ausgabe
    \texttt{ir\_output.json} als Python-Dictionary (\texttt{ir\_data}).
    \item \textbf{Trockenaudio laden:} \texttt{read\_wav} liest die
    Eingangsdatei und liefert Audiodaten samt Abtastrate.
    \item \textbf{Faltung:} \texttt{convolve} wendet die aus \texttt{ir\_data}
    erzeugte Impulsantwort auf das Trockensignal an.
    \item \textbf{Export:} \texttt{write\_wav} normalisiert das Ergebnis bei
    Bedarf und schreibt die fertige Ausgabedatei.
\end{enumerate}

Der Binary-Name wird plattformabhängig gewählt: unter Windows mit der Endung
\texttt{.exe}, unter macOS und Linux ohne. Stimmt die Abtastrate des Trockenaudios
nicht mit der Rate der Engine überein, gibt \texttt{main.py} eine Warnung aus,
weil der Hall sonst verstimmt klingen kann.

\subsection{Die Prozessbrücke zwischen Rust und Python}
\label{subsec:rust-runner}

Rust und Python sind nicht über eine Sprachbindung wie PyO3 gekoppelt, sondern
lose über einen Subprozess-Aufruf und eine JSON-Datei verbunden. Die Funktion
\texttt{run\_rust\_and\_load\_json} kapselt diese Brücke. Sie prüft zunächst, ob
die Binary existiert, und bestimmt das Arbeitsverzeichnis anhand der
\texttt{input\_config.json}. Eine alte Ausgabedatei wird vorher gelöscht, damit
kein veraltetes Ergebnis eingelesen wird. Nach dem Lauf wertet die Funktion den
Rückgabecode aus; schlägt Rust fehl, landen \texttt{stdout} und \texttt{stderr} in
der Fehlermeldung und erleichtern die Diagnose. Der Vorteil der losen Kopplung
liegt im verteilten Betrieb: Beide Dienste lassen sich getrennt kompilieren,
testen und auf mehrere Cluster-Knoten verteilen.

\subsection{Signalverarbeitung mit SciPy}
\label{subsec:scipy}

Das Modul \texttt{convolution.py} bildet den Kern des Python-Dienstes. Es wandelt
die rohen Verzögerungs- und Druckdaten in eine abtastgenaue Impulsantwort um und
faltet damit das Trockensignal.

\subsubsection{Diskretisierung der Impulsantwort}
\label{subsubsec:build-ir}

Die Rust-Engine liefert eine endliche Menge diskreter Strahl-Treffer, jeweils
bestehend aus einer Ankunftsverzögerung $\tau_j$ (in Sekunden) und einem skalaren
Restdruck $p_j$. Die Funktion \texttt{build\_impulse\_response} quantisiert jeden
Treffer auf den nächstgelegenen Abtastindex $i_j = \operatorname{round}(\tau_j
\cdot f_s)$ und akkumuliert gleichzeitige Treffer zur diskreten Impulsantwort:
\[
    h[i] \;=\; \sum_{j} p_j \cdot \bigl[\,i_j = i\,\bigr].
\]
Die Akkumulation erfolgt mit \texttt{numpy.add.at}. Anders als die einfache
indizierte Zuweisung, die bei doppelten Indizes nur den letzten Wert behält,
summiert \texttt{add.at} gleichzeitige Treffer korrekt auf. Vor der Berechnung
prüft die Funktion die Eingaben an der Systemgrenze: übereinstimmende
Array-Formen, Eindimensionalität, nicht-negative Verzögerungen und eine positive
Abtastrate.

\begin{lstlisting}[style=cemdsp, caption={Diskretisierung der Ray-Hits zur Impulsantwort.}]
indices = np.round(delays * sample_rate).astype(np.int64)
ir_length = int(indices.max()) + 1
ir = np.zeros(ir_length, dtype=np.float64)
np.add.at(ir, indices, pressures)   # akkumuliert gleichzeitige Treffer
return ir
\end{lstlisting}

\subsubsection{Faltung mit \texttt{oaconvolve}}
\label{subsubsec:oaconvolve}

Die eigentliche Faltung nutzt \texttt{scipy.signal.oaconvolve}, die
Overlap-Add-basierte FFT-Faltung aus Abschnitt~2.3. Der Modus \texttt{"full"}
liefert das vollständige Ergebnis der Länge $N + M - 1$ und bewahrt damit den
kompletten Nachhall. Mehrkanaliges (etwa Stereo-)Audio wird kanalweise gefaltet
und danach wieder zusammengesetzt:

\begin{lstlisting}[style=cemdsp, caption={Mono- und Mehrkanal-Faltung.}]
if dry_audio.ndim == 1:
    wet = signal.oaconvolve(dry_audio, ir, mode="full")
elif dry_audio.ndim == 2:
    channels = [
        signal.oaconvolve(dry_audio[:, c], ir, mode="full")
        for c in range(dry_audio.shape[1])
    ]
    wet = np.stack(channels, axis=1)
return wet.astype(np.float32)
\end{lstlisting}

Vor der Faltung prüft \texttt{convolve} das Schema von \texttt{ir\_data} und
vergleicht die Abtastraten von Audio und Impulsantwort. Das Ergebnis gibt die
Funktion roh und unnormalisiert zurück. Die Lautstärke regelt erst der Export
(Abschnitt~\ref{subsec:rendering}); so bleibt die Faltung frei von Pegelfragen.

\subsubsection{Frequenzabhängige Materialdämpfung}
\label{subsubsec:material}

Reale Materialien dämpfen Schall frequenzabhängig. Ein Teppich schluckt hohe
Frequenzen viel stärker als tiefe, Beton reflektiert dagegen fast breitbandig.
Die Rust-Engine reduziert die Absorption aber schon auf einen einzigen skalaren
Restdruck $p_j$ pro Treffer, sodass die spektrale Information an dieser
Systemgrenze verloren geht. Damit der Nachhall trotzdem frequenzabhängig wird,
ergänzt \texttt{convolution.py} ein Absorptionsmodell im Frequenzbereich.

Grundlage sind die Oktavband-Absorptionskoeffizienten $\alpha_b$ aus der Datei
\texttt{materials.json}, die auch die Rust-Engine speist. Physik und
Signalverarbeitung greifen damit auf dieselbe Datenbasis zu (single source of
truth). Jedes Material ist über sechs Oktavbänder mit den Mittenfrequenzen
$125,\,250,\,500,\,1000,\,2000$ und $4000$~Hz beschrieben; ein hoher Koeffizient
bedeutet starke Dämpfung im jeweiligen Band.

Das Modell beruht auf einer einfachen Überlegung: Ein Tap, der zum Zeitpunkt $t$
eintrifft, hat näherungsweise $n(t) = t / \tau_{\mathrm{mfp}}$ Reflexionen hinter
sich, wobei $\tau_{\mathrm{mfp}}$ die mittlere Zeit zwischen zwei Reflexionen
angibt (die mittlere freie Weglänge in Zeiteinheiten). Nach $n$ Reflexionen
verbleibt in jedem Oktavband $b$ der Energieanteil
\[
    g_b(t) \;=\; \bigl(1 - \alpha_b\bigr)^{\,n(t)} \;=\; \bigl(1 - \alpha_b\bigr)^{\,t / \tau_{\mathrm{mfp}}}.
\]
Bänder mit hohem $\alpha_b$, etwa die Höhen bei Teppich, klingen dadurch deutlich
schneller ab als solche mit niedrigem $\alpha_b$. Der Nachhall wird über die Zeit
zunehmend dunkler.

Im Code setzt eine Oktavband-Filterbank dieses Modell um. Die Impulsantwort wird
über Butterworth-Bandpässe (nullphasig via \texttt{sosfiltfilt}) in ihre sechs
Bänder zerlegt, jedes Band mit der zeitabhängigen Hüllkurve $g_b(t)$ gewichtet und
die Bänder anschließend wieder aufsummiert.

\begin{lstlisting}[style=cemdsp, caption={Frequenzabhaengige Material-Daempfung der Impulsantwort.}]
centers, materials = _load_materials(str(_MATERIALS_JSON))
nyquist = sample_rate / 2.0
reflections = (np.arange(ir.size) / sample_rate) / mean_free_path_s

shaped = np.zeros_like(ir)
for (low, high), alpha in zip(_octave_band_edges(centers), materials[material]):
    band = _bandpass(ir, low, high, nyquist)      # ein Oktavband isolieren
    gain = np.power(max(1e-6, 1.0 - alpha), reflections)
    shaped += band * gain                          # zeitabhaengig daempfen
return shaped
\end{lstlisting}

\paragraph{Verifikation.} Ein reproduzierbarer Selbsttest prüft das Modell. Er
dämpft dieselbe dichte Impulsantwort einmal mit \texttt{carpet} und einmal mit
\texttt{concrete} und misst jeweils den Energieanteil oberhalb von $2$~kHz. Bei
Teppich bleiben davon nur $0{,}39$ übrig, bei Beton $0{,}71$; der Teppich schluckt
die Höhen also deutlich stärker. Kippt diese Reihenfolge, schlägt der Test fehl.

\paragraph{Grenzen des Modells.} Das Verfahren modelliert ein globales
Raummaterial, nicht die tatsächliche Materialfolge eines einzelnen Strahlpfads.
Eine Trennung nach Oberfläche, also die Aussage „dieser Tap stammt von Teppich,
jener von Beton", ist im Python-Dienst nicht möglich, solange die Rust-Engine nur
einen skalaren Druck pro Treffer liefert. Dafür müsste die Engine schon bandweise
Drücke ($p_{j,b}$ statt $p_j$) exportieren. Kapitel~7 führt das als mögliche
Erweiterung.

\subsection{Audio-Rendering und Export}
\label{subsec:rendering}

Das Ein- und Auslesen der Audiodateien sowie die Normalisierung gegen Clipping
übernimmt \texttt{wav\_handler.py}; die Details stehen in der Teildokumentation zum
Audio-I/O. Für den DSP-Teil zählt nur der Vertrag zwischen den Stufen.
\texttt{convolve} liefert ein rohes, womöglich übersteuertes Signal;
\texttt{write\_wav} misst dessen maximale Amplitude und normalisiert erst, wenn
sie den Bereich $[-1{,}0,\,1{,}0]$ verlässt. So bleibt Clipping aus, und die
Faltung muss sich nie um Pegel kümmern.

%   3. Der Master braucht nur diese Pakete (Dokumentation.tex hat beide schon):
%        \usepackage{amsmath}
%        \usepackage{listings}
%      Kein xcolor/babel noetig: Umlaute laufen unter UTF-8, deutsche
%      Anfuehrungszeichen sind literal „..." gesetzt.
%
% STANDALONE-VORSCHAU (zum alleine Kompilieren einfach drumherum setzen):
%   \documentclass{article}\usepackage{amsmath}\usepackage{listings}
%   \begin{document}
%   % =====================================================================
% Kapitel 5 - Digitale Signalverarbeitung und Orchestrierung (Python)
% Autor: Cem Acar
%
% EINBINDEN ins Hauptdokument:
%   1. Dieses File nach latexDocs/ legen.
%   2. Im Master (Dokumentation.tex) an der richtigen Stelle einfuegen:
%        \input{Kapitel5_DSP_Cem}
%   3. Der Master braucht nur diese Pakete (Dokumentation.tex hat beide schon):
%        \usepackage{amsmath}
%        \usepackage{listings}
%      Kein xcolor/babel noetig: Umlaute laufen unter UTF-8, deutsche
%      Anfuehrungszeichen sind literal „..." gesetzt.
%
% STANDALONE-VORSCHAU (zum alleine Kompilieren einfach drumherum setzen):
%   \documentclass{article}\usepackage{amsmath}\usepackage{listings}
%   \begin{document}
%   \input{Kapitel5_DSP_Cem}
%   \end{document}
% =====================================================================

% Listing-Stil fuer dieses Kapitel (nur listings-Kern, keine Farben noetig)
\lstdefinestyle{cemdsp}{
    language=Python,
    basicstyle=\ttfamily\small,
    showstringspaces=false,
    breaklines=true,
    numbers=left,
    numberstyle=\tiny,
    frame=single,
    tabsize=4
}

%======================================================================
\section{Digitale Signalverarbeitung und Orchestrierung (Python)}
\label{sec:python-dsp}

Die Rust-Engine berechnet die Physik der Schallausbreitung, der Python-Dienst
übernimmt die Signalverarbeitung und die Orchestrierung der Pipeline. Warum die
Physik in Rust und die Signalverarbeitung in Python mit NumPy und SciPy liegt,
ist in Kapitel~3 begründet. Die signaltheoretischen Grundlagen, also
Impulsantwort, Faltungsintegral und schnelle Faltung über die FFT, behandelt
Abschnitt~2.3. Dieses Kapitel zeigt, wie diese Grundlagen im Code umgesetzt sind,
und folgt dabei dem Datenfluss durch die Pipeline.

\subsection{Pipeline-Management durch \texttt{main.py}}
\label{subsec:pipeline}

\texttt{main.py} führt die vier Verarbeitungsschritte nacheinander aus. Tritt in
einem Schritt ein Fehler auf, bricht die Pipeline sofort mit einer Meldung ab,
statt mit ungültigen Zwischendaten weiterzurechnen (fail-fast).

\begin{enumerate}
    \item \textbf{Physik-Engine starten:} \texttt{run\_rust\_and\_load\_json} ruft
    die vorkompilierte Rust-Binary als Subprozess auf und lädt deren Ausgabe
    \texttt{ir\_output.json} als Python-Dictionary (\texttt{ir\_data}).
    \item \textbf{Trockenaudio laden:} \texttt{read\_wav} liest die
    Eingangsdatei und liefert Audiodaten samt Abtastrate.
    \item \textbf{Faltung:} \texttt{convolve} wendet die aus \texttt{ir\_data}
    erzeugte Impulsantwort auf das Trockensignal an.
    \item \textbf{Export:} \texttt{write\_wav} normalisiert das Ergebnis bei
    Bedarf und schreibt die fertige Ausgabedatei.
\end{enumerate}

Der Binary-Name wird plattformabhängig gewählt: unter Windows mit der Endung
\texttt{.exe}, unter macOS und Linux ohne. Stimmt die Abtastrate des Trockenaudios
nicht mit der Rate der Engine überein, gibt \texttt{main.py} eine Warnung aus,
weil der Hall sonst verstimmt klingen kann.

\subsection{Die Prozessbrücke zwischen Rust und Python}
\label{subsec:rust-runner}

Rust und Python sind nicht über eine Sprachbindung wie PyO3 gekoppelt, sondern
lose über einen Subprozess-Aufruf und eine JSON-Datei verbunden. Die Funktion
\texttt{run\_rust\_and\_load\_json} kapselt diese Brücke. Sie prüft zunächst, ob
die Binary existiert, und bestimmt das Arbeitsverzeichnis anhand der
\texttt{input\_config.json}. Eine alte Ausgabedatei wird vorher gelöscht, damit
kein veraltetes Ergebnis eingelesen wird. Nach dem Lauf wertet die Funktion den
Rückgabecode aus; schlägt Rust fehl, landen \texttt{stdout} und \texttt{stderr} in
der Fehlermeldung und erleichtern die Diagnose. Der Vorteil der losen Kopplung
liegt im verteilten Betrieb: Beide Dienste lassen sich getrennt kompilieren,
testen und auf mehrere Cluster-Knoten verteilen.

\subsection{Signalverarbeitung mit SciPy}
\label{subsec:scipy}

Das Modul \texttt{convolution.py} bildet den Kern des Python-Dienstes. Es wandelt
die rohen Verzögerungs- und Druckdaten in eine abtastgenaue Impulsantwort um und
faltet damit das Trockensignal.

\subsubsection{Diskretisierung der Impulsantwort}
\label{subsubsec:build-ir}

Die Rust-Engine liefert eine endliche Menge diskreter Strahl-Treffer, jeweils
bestehend aus einer Ankunftsverzögerung $\tau_j$ (in Sekunden) und einem skalaren
Restdruck $p_j$. Die Funktion \texttt{build\_impulse\_response} quantisiert jeden
Treffer auf den nächstgelegenen Abtastindex $i_j = \operatorname{round}(\tau_j
\cdot f_s)$ und akkumuliert gleichzeitige Treffer zur diskreten Impulsantwort:
\[
    h[i] \;=\; \sum_{j} p_j \cdot \bigl[\,i_j = i\,\bigr].
\]
Die Akkumulation erfolgt mit \texttt{numpy.add.at}. Anders als die einfache
indizierte Zuweisung, die bei doppelten Indizes nur den letzten Wert behält,
summiert \texttt{add.at} gleichzeitige Treffer korrekt auf. Vor der Berechnung
prüft die Funktion die Eingaben an der Systemgrenze: übereinstimmende
Array-Formen, Eindimensionalität, nicht-negative Verzögerungen und eine positive
Abtastrate.

\begin{lstlisting}[style=cemdsp, caption={Diskretisierung der Ray-Hits zur Impulsantwort.}]
indices = np.round(delays * sample_rate).astype(np.int64)
ir_length = int(indices.max()) + 1
ir = np.zeros(ir_length, dtype=np.float64)
np.add.at(ir, indices, pressures)   # akkumuliert gleichzeitige Treffer
return ir
\end{lstlisting}

\subsubsection{Faltung mit \texttt{oaconvolve}}
\label{subsubsec:oaconvolve}

Die eigentliche Faltung nutzt \texttt{scipy.signal.oaconvolve}, die
Overlap-Add-basierte FFT-Faltung aus Abschnitt~2.3. Der Modus \texttt{"full"}
liefert das vollständige Ergebnis der Länge $N + M - 1$ und bewahrt damit den
kompletten Nachhall. Mehrkanaliges (etwa Stereo-)Audio wird kanalweise gefaltet
und danach wieder zusammengesetzt:

\begin{lstlisting}[style=cemdsp, caption={Mono- und Mehrkanal-Faltung.}]
if dry_audio.ndim == 1:
    wet = signal.oaconvolve(dry_audio, ir, mode="full")
elif dry_audio.ndim == 2:
    channels = [
        signal.oaconvolve(dry_audio[:, c], ir, mode="full")
        for c in range(dry_audio.shape[1])
    ]
    wet = np.stack(channels, axis=1)
return wet.astype(np.float32)
\end{lstlisting}

Vor der Faltung prüft \texttt{convolve} das Schema von \texttt{ir\_data} und
vergleicht die Abtastraten von Audio und Impulsantwort. Das Ergebnis gibt die
Funktion roh und unnormalisiert zurück. Die Lautstärke regelt erst der Export
(Abschnitt~\ref{subsec:rendering}); so bleibt die Faltung frei von Pegelfragen.

\subsubsection{Frequenzabhängige Materialdämpfung}
\label{subsubsec:material}

Reale Materialien dämpfen Schall frequenzabhängig. Ein Teppich schluckt hohe
Frequenzen viel stärker als tiefe, Beton reflektiert dagegen fast breitbandig.
Die Rust-Engine reduziert die Absorption aber schon auf einen einzigen skalaren
Restdruck $p_j$ pro Treffer, sodass die spektrale Information an dieser
Systemgrenze verloren geht. Damit der Nachhall trotzdem frequenzabhängig wird,
ergänzt \texttt{convolution.py} ein Absorptionsmodell im Frequenzbereich.

Grundlage sind die Oktavband-Absorptionskoeffizienten $\alpha_b$ aus der Datei
\texttt{materials.json}, die auch die Rust-Engine speist. Physik und
Signalverarbeitung greifen damit auf dieselbe Datenbasis zu (single source of
truth). Jedes Material ist über sechs Oktavbänder mit den Mittenfrequenzen
$125,\,250,\,500,\,1000,\,2000$ und $4000$~Hz beschrieben; ein hoher Koeffizient
bedeutet starke Dämpfung im jeweiligen Band.

Das Modell beruht auf einer einfachen Überlegung: Ein Tap, der zum Zeitpunkt $t$
eintrifft, hat näherungsweise $n(t) = t / \tau_{\mathrm{mfp}}$ Reflexionen hinter
sich, wobei $\tau_{\mathrm{mfp}}$ die mittlere Zeit zwischen zwei Reflexionen
angibt (die mittlere freie Weglänge in Zeiteinheiten). Nach $n$ Reflexionen
verbleibt in jedem Oktavband $b$ der Energieanteil
\[
    g_b(t) \;=\; \bigl(1 - \alpha_b\bigr)^{\,n(t)} \;=\; \bigl(1 - \alpha_b\bigr)^{\,t / \tau_{\mathrm{mfp}}}.
\]
Bänder mit hohem $\alpha_b$, etwa die Höhen bei Teppich, klingen dadurch deutlich
schneller ab als solche mit niedrigem $\alpha_b$. Der Nachhall wird über die Zeit
zunehmend dunkler.

Im Code setzt eine Oktavband-Filterbank dieses Modell um. Die Impulsantwort wird
über Butterworth-Bandpässe (nullphasig via \texttt{sosfiltfilt}) in ihre sechs
Bänder zerlegt, jedes Band mit der zeitabhängigen Hüllkurve $g_b(t)$ gewichtet und
die Bänder anschließend wieder aufsummiert.

\begin{lstlisting}[style=cemdsp, caption={Frequenzabhaengige Material-Daempfung der Impulsantwort.}]
centers, materials = _load_materials(str(_MATERIALS_JSON))
nyquist = sample_rate / 2.0
reflections = (np.arange(ir.size) / sample_rate) / mean_free_path_s

shaped = np.zeros_like(ir)
for (low, high), alpha in zip(_octave_band_edges(centers), materials[material]):
    band = _bandpass(ir, low, high, nyquist)      # ein Oktavband isolieren
    gain = np.power(max(1e-6, 1.0 - alpha), reflections)
    shaped += band * gain                          # zeitabhaengig daempfen
return shaped
\end{lstlisting}

\paragraph{Verifikation.} Ein reproduzierbarer Selbsttest prüft das Modell. Er
dämpft dieselbe dichte Impulsantwort einmal mit \texttt{carpet} und einmal mit
\texttt{concrete} und misst jeweils den Energieanteil oberhalb von $2$~kHz. Bei
Teppich bleiben davon nur $0{,}39$ übrig, bei Beton $0{,}71$; der Teppich schluckt
die Höhen also deutlich stärker. Kippt diese Reihenfolge, schlägt der Test fehl.

\paragraph{Grenzen des Modells.} Das Verfahren modelliert ein globales
Raummaterial, nicht die tatsächliche Materialfolge eines einzelnen Strahlpfads.
Eine Trennung nach Oberfläche, also die Aussage „dieser Tap stammt von Teppich,
jener von Beton", ist im Python-Dienst nicht möglich, solange die Rust-Engine nur
einen skalaren Druck pro Treffer liefert. Dafür müsste die Engine schon bandweise
Drücke ($p_{j,b}$ statt $p_j$) exportieren. Kapitel~7 führt das als mögliche
Erweiterung.

\subsection{Audio-Rendering und Export}
\label{subsec:rendering}

Das Ein- und Auslesen der Audiodateien sowie die Normalisierung gegen Clipping
übernimmt \texttt{wav\_handler.py}; die Details stehen in der Teildokumentation zum
Audio-I/O. Für den DSP-Teil zählt nur der Vertrag zwischen den Stufen.
\texttt{convolve} liefert ein rohes, womöglich übersteuertes Signal;
\texttt{write\_wav} misst dessen maximale Amplitude und normalisiert erst, wenn
sie den Bereich $[-1{,}0,\,1{,}0]$ verlässt. So bleibt Clipping aus, und die
Faltung muss sich nie um Pegel kümmern.

%   \end{document}
% =====================================================================

% Listing-Stil fuer dieses Kapitel (nur listings-Kern, keine Farben noetig)
\lstdefinestyle{cemdsp}{
    language=Python,
    basicstyle=\ttfamily\small,
    showstringspaces=false,
    breaklines=true,
    numbers=left,
    numberstyle=\tiny,
    frame=single,
    tabsize=4
}

%======================================================================
\section{Digitale Signalverarbeitung und Orchestrierung (Python)}
\label{sec:python-dsp}

Die Rust-Engine berechnet die Physik der Schallausbreitung, der Python-Dienst
übernimmt die Signalverarbeitung und die Orchestrierung der Pipeline. Warum die
Physik in Rust und die Signalverarbeitung in Python mit NumPy und SciPy liegt,
ist in Kapitel~3 begründet. Die signaltheoretischen Grundlagen, also
Impulsantwort, Faltungsintegral und schnelle Faltung über die FFT, behandelt
Abschnitt~2.3. Dieses Kapitel zeigt, wie diese Grundlagen im Code umgesetzt sind,
und folgt dabei dem Datenfluss durch die Pipeline.

\subsection{Pipeline-Management durch \texttt{main.py}}
\label{subsec:pipeline}

\texttt{main.py} führt die vier Verarbeitungsschritte nacheinander aus. Tritt in
einem Schritt ein Fehler auf, bricht die Pipeline sofort mit einer Meldung ab,
statt mit ungültigen Zwischendaten weiterzurechnen (fail-fast).

\begin{enumerate}
    \item \textbf{Physik-Engine starten:} \texttt{run\_rust\_and\_load\_json} ruft
    die vorkompilierte Rust-Binary als Subprozess auf und lädt deren Ausgabe
    \texttt{ir\_output.json} als Python-Dictionary (\texttt{ir\_data}).
    \item \textbf{Trockenaudio laden:} \texttt{read\_wav} liest die
    Eingangsdatei und liefert Audiodaten samt Abtastrate.
    \item \textbf{Faltung:} \texttt{convolve} wendet die aus \texttt{ir\_data}
    erzeugte Impulsantwort auf das Trockensignal an.
    \item \textbf{Export:} \texttt{write\_wav} normalisiert das Ergebnis bei
    Bedarf und schreibt die fertige Ausgabedatei.
\end{enumerate}

Der Binary-Name wird plattformabhängig gewählt: unter Windows mit der Endung
\texttt{.exe}, unter macOS und Linux ohne. Stimmt die Abtastrate des Trockenaudios
nicht mit der Rate der Engine überein, gibt \texttt{main.py} eine Warnung aus,
weil der Hall sonst verstimmt klingen kann.

\subsection{Die Prozessbrücke zwischen Rust und Python}
\label{subsec:rust-runner}

Rust und Python sind nicht über eine Sprachbindung wie PyO3 gekoppelt, sondern
lose über einen Subprozess-Aufruf und eine JSON-Datei verbunden. Die Funktion
\texttt{run\_rust\_and\_load\_json} kapselt diese Brücke. Sie prüft zunächst, ob
die Binary existiert, und bestimmt das Arbeitsverzeichnis anhand der
\texttt{input\_config.json}. Eine alte Ausgabedatei wird vorher gelöscht, damit
kein veraltetes Ergebnis eingelesen wird. Nach dem Lauf wertet die Funktion den
Rückgabecode aus; schlägt Rust fehl, landen \texttt{stdout} und \texttt{stderr} in
der Fehlermeldung und erleichtern die Diagnose. Der Vorteil der losen Kopplung
liegt im verteilten Betrieb: Beide Dienste lassen sich getrennt kompilieren,
testen und auf mehrere Cluster-Knoten verteilen.

\subsection{Signalverarbeitung mit SciPy}
\label{subsec:scipy}

Das Modul \texttt{convolution.py} bildet den Kern des Python-Dienstes. Es wandelt
die rohen Verzögerungs- und Druckdaten in eine abtastgenaue Impulsantwort um und
faltet damit das Trockensignal.

\subsubsection{Diskretisierung der Impulsantwort}
\label{subsubsec:build-ir}

Die Rust-Engine liefert eine endliche Menge diskreter Strahl-Treffer, jeweils
bestehend aus einer Ankunftsverzögerung $\tau_j$ (in Sekunden) und einem skalaren
Restdruck $p_j$. Die Funktion \texttt{build\_impulse\_response} quantisiert jeden
Treffer auf den nächstgelegenen Abtastindex $i_j = \operatorname{round}(\tau_j
\cdot f_s)$ und akkumuliert gleichzeitige Treffer zur diskreten Impulsantwort:
\[
    h[i] \;=\; \sum_{j} p_j \cdot \bigl[\,i_j = i\,\bigr].
\]
Die Akkumulation erfolgt mit \texttt{numpy.add.at}. Anders als die einfache
indizierte Zuweisung, die bei doppelten Indizes nur den letzten Wert behält,
summiert \texttt{add.at} gleichzeitige Treffer korrekt auf. Vor der Berechnung
prüft die Funktion die Eingaben an der Systemgrenze: übereinstimmende
Array-Formen, Eindimensionalität, nicht-negative Verzögerungen und eine positive
Abtastrate.

\begin{lstlisting}[style=cemdsp, caption={Diskretisierung der Ray-Hits zur Impulsantwort.}]
indices = np.round(delays * sample_rate).astype(np.int64)
ir_length = int(indices.max()) + 1
ir = np.zeros(ir_length, dtype=np.float64)
np.add.at(ir, indices, pressures)   # akkumuliert gleichzeitige Treffer
return ir
\end{lstlisting}

\subsubsection{Faltung mit \texttt{oaconvolve}}
\label{subsubsec:oaconvolve}

Die eigentliche Faltung nutzt \texttt{scipy.signal.oaconvolve}, die
Overlap-Add-basierte FFT-Faltung aus Abschnitt~2.3. Der Modus \texttt{"full"}
liefert das vollständige Ergebnis der Länge $N + M - 1$ und bewahrt damit den
kompletten Nachhall. Mehrkanaliges (etwa Stereo-)Audio wird kanalweise gefaltet
und danach wieder zusammengesetzt:

\begin{lstlisting}[style=cemdsp, caption={Mono- und Mehrkanal-Faltung.}]
if dry_audio.ndim == 1:
    wet = signal.oaconvolve(dry_audio, ir, mode="full")
elif dry_audio.ndim == 2:
    channels = [
        signal.oaconvolve(dry_audio[:, c], ir, mode="full")
        for c in range(dry_audio.shape[1])
    ]
    wet = np.stack(channels, axis=1)
return wet.astype(np.float32)
\end{lstlisting}

Vor der Faltung prüft \texttt{convolve} das Schema von \texttt{ir\_data} und
vergleicht die Abtastraten von Audio und Impulsantwort. Das Ergebnis gibt die
Funktion roh und unnormalisiert zurück. Die Lautstärke regelt erst der Export
(Abschnitt~\ref{subsec:rendering}); so bleibt die Faltung frei von Pegelfragen.

\subsubsection{Frequenzabhängige Materialdämpfung}
\label{subsubsec:material}

Reale Materialien dämpfen Schall frequenzabhängig. Ein Teppich schluckt hohe
Frequenzen viel stärker als tiefe, Beton reflektiert dagegen fast breitbandig.
Die Rust-Engine reduziert die Absorption aber schon auf einen einzigen skalaren
Restdruck $p_j$ pro Treffer, sodass die spektrale Information an dieser
Systemgrenze verloren geht. Damit der Nachhall trotzdem frequenzabhängig wird,
ergänzt \texttt{convolution.py} ein Absorptionsmodell im Frequenzbereich.

Grundlage sind die Oktavband-Absorptionskoeffizienten $\alpha_b$ aus der Datei
\texttt{materials.json}, die auch die Rust-Engine speist. Physik und
Signalverarbeitung greifen damit auf dieselbe Datenbasis zu (single source of
truth). Jedes Material ist über sechs Oktavbänder mit den Mittenfrequenzen
$125,\,250,\,500,\,1000,\,2000$ und $4000$~Hz beschrieben; ein hoher Koeffizient
bedeutet starke Dämpfung im jeweiligen Band.

Das Modell beruht auf einer einfachen Überlegung: Ein Tap, der zum Zeitpunkt $t$
eintrifft, hat näherungsweise $n(t) = t / \tau_{\mathrm{mfp}}$ Reflexionen hinter
sich, wobei $\tau_{\mathrm{mfp}}$ die mittlere Zeit zwischen zwei Reflexionen
angibt (die mittlere freie Weglänge in Zeiteinheiten). Nach $n$ Reflexionen
verbleibt in jedem Oktavband $b$ der Energieanteil
\[
    g_b(t) \;=\; \bigl(1 - \alpha_b\bigr)^{\,n(t)} \;=\; \bigl(1 - \alpha_b\bigr)^{\,t / \tau_{\mathrm{mfp}}}.
\]
Bänder mit hohem $\alpha_b$, etwa die Höhen bei Teppich, klingen dadurch deutlich
schneller ab als solche mit niedrigem $\alpha_b$. Der Nachhall wird über die Zeit
zunehmend dunkler.

Im Code setzt eine Oktavband-Filterbank dieses Modell um. Die Impulsantwort wird
über Butterworth-Bandpässe (nullphasig via \texttt{sosfiltfilt}) in ihre sechs
Bänder zerlegt, jedes Band mit der zeitabhängigen Hüllkurve $g_b(t)$ gewichtet und
die Bänder anschließend wieder aufsummiert.

\begin{lstlisting}[style=cemdsp, caption={Frequenzabhaengige Material-Daempfung der Impulsantwort.}]
centers, materials = _load_materials(str(_MATERIALS_JSON))
nyquist = sample_rate / 2.0
reflections = (np.arange(ir.size) / sample_rate) / mean_free_path_s

shaped = np.zeros_like(ir)
for (low, high), alpha in zip(_octave_band_edges(centers), materials[material]):
    band = _bandpass(ir, low, high, nyquist)      # ein Oktavband isolieren
    gain = np.power(max(1e-6, 1.0 - alpha), reflections)
    shaped += band * gain                          # zeitabhaengig daempfen
return shaped
\end{lstlisting}

\paragraph{Verifikation.} Ein reproduzierbarer Selbsttest prüft das Modell. Er
dämpft dieselbe dichte Impulsantwort einmal mit \texttt{carpet} und einmal mit
\texttt{concrete} und misst jeweils den Energieanteil oberhalb von $2$~kHz. Bei
Teppich bleiben davon nur $0{,}39$ übrig, bei Beton $0{,}71$; der Teppich schluckt
die Höhen also deutlich stärker. Kippt diese Reihenfolge, schlägt der Test fehl.

\paragraph{Grenzen des Modells.} Das Verfahren modelliert ein globales
Raummaterial, nicht die tatsächliche Materialfolge eines einzelnen Strahlpfads.
Eine Trennung nach Oberfläche, also die Aussage „dieser Tap stammt von Teppich,
jener von Beton", ist im Python-Dienst nicht möglich, solange die Rust-Engine nur
einen skalaren Druck pro Treffer liefert. Dafür müsste die Engine schon bandweise
Drücke ($p_{j,b}$ statt $p_j$) exportieren. Kapitel~7 führt das als mögliche
Erweiterung.

\subsection{Audio-Rendering und Export}
\label{subsec:rendering}

Das Ein- und Auslesen der Audiodateien sowie die Normalisierung gegen Clipping
übernimmt \texttt{wav\_handler.py}; die Details stehen in der Teildokumentation zum
Audio-I/O. Für den DSP-Teil zählt nur der Vertrag zwischen den Stufen.
\texttt{convolve} liefert ein rohes, womöglich übersteuertes Signal;
\texttt{write\_wav} misst dessen maximale Amplitude und normalisiert erst, wenn
sie den Bereich $[-1{,}0,\,1{,}0]$ verlässt. So bleibt Clipping aus, und die
Faltung muss sich nie um Pegel kümmern.

%   3. Der Master braucht nur diese Pakete (Dokumentation.tex hat beide schon):
%        \usepackage{amsmath}
%        \usepackage{listings}
%      Kein xcolor/babel noetig: Umlaute laufen unter UTF-8, deutsche
%      Anfuehrungszeichen sind literal „..." gesetzt.
%
% STANDALONE-VORSCHAU (zum alleine Kompilieren einfach drumherum setzen):
%   \documentclass{article}\usepackage{amsmath}\usepackage{listings}
%   \begin{document}
%   % =====================================================================
% Kapitel 5 - Digitale Signalverarbeitung und Orchestrierung (Python)
% Autor: Cem Acar
%
% EINBINDEN ins Hauptdokument:
%   1. Dieses File nach latexDocs/ legen.
%   2. Im Master (Dokumentation.tex) an der richtigen Stelle einfuegen:
%        % =====================================================================
% Kapitel 5 - Digitale Signalverarbeitung und Orchestrierung (Python)
% Autor: Cem Acar
%
% EINBINDEN ins Hauptdokument:
%   1. Dieses File nach latexDocs/ legen.
%   2. Im Master (Dokumentation.tex) an der richtigen Stelle einfuegen:
%        \input{Kapitel5_DSP_Cem}
%   3. Der Master braucht nur diese Pakete (Dokumentation.tex hat beide schon):
%        \usepackage{amsmath}
%        \usepackage{listings}
%      Kein xcolor/babel noetig: Umlaute laufen unter UTF-8, deutsche
%      Anfuehrungszeichen sind literal „..." gesetzt.
%
% STANDALONE-VORSCHAU (zum alleine Kompilieren einfach drumherum setzen):
%   \documentclass{article}\usepackage{amsmath}\usepackage{listings}
%   \begin{document}
%   \input{Kapitel5_DSP_Cem}
%   \end{document}
% =====================================================================

% Listing-Stil fuer dieses Kapitel (nur listings-Kern, keine Farben noetig)
\lstdefinestyle{cemdsp}{
    language=Python,
    basicstyle=\ttfamily\small,
    showstringspaces=false,
    breaklines=true,
    numbers=left,
    numberstyle=\tiny,
    frame=single,
    tabsize=4
}

%======================================================================
\section{Digitale Signalverarbeitung und Orchestrierung (Python)}
\label{sec:python-dsp}

Die Rust-Engine berechnet die Physik der Schallausbreitung, der Python-Dienst
übernimmt die Signalverarbeitung und die Orchestrierung der Pipeline. Warum die
Physik in Rust und die Signalverarbeitung in Python mit NumPy und SciPy liegt,
ist in Kapitel~3 begründet. Die signaltheoretischen Grundlagen, also
Impulsantwort, Faltungsintegral und schnelle Faltung über die FFT, behandelt
Abschnitt~2.3. Dieses Kapitel zeigt, wie diese Grundlagen im Code umgesetzt sind,
und folgt dabei dem Datenfluss durch die Pipeline.

\subsection{Pipeline-Management durch \texttt{main.py}}
\label{subsec:pipeline}

\texttt{main.py} führt die vier Verarbeitungsschritte nacheinander aus. Tritt in
einem Schritt ein Fehler auf, bricht die Pipeline sofort mit einer Meldung ab,
statt mit ungültigen Zwischendaten weiterzurechnen (fail-fast).

\begin{enumerate}
    \item \textbf{Physik-Engine starten:} \texttt{run\_rust\_and\_load\_json} ruft
    die vorkompilierte Rust-Binary als Subprozess auf und lädt deren Ausgabe
    \texttt{ir\_output.json} als Python-Dictionary (\texttt{ir\_data}).
    \item \textbf{Trockenaudio laden:} \texttt{read\_wav} liest die
    Eingangsdatei und liefert Audiodaten samt Abtastrate.
    \item \textbf{Faltung:} \texttt{convolve} wendet die aus \texttt{ir\_data}
    erzeugte Impulsantwort auf das Trockensignal an.
    \item \textbf{Export:} \texttt{write\_wav} normalisiert das Ergebnis bei
    Bedarf und schreibt die fertige Ausgabedatei.
\end{enumerate}

Der Binary-Name wird plattformabhängig gewählt: unter Windows mit der Endung
\texttt{.exe}, unter macOS und Linux ohne. Stimmt die Abtastrate des Trockenaudios
nicht mit der Rate der Engine überein, gibt \texttt{main.py} eine Warnung aus,
weil der Hall sonst verstimmt klingen kann.

\subsection{Die Prozessbrücke zwischen Rust und Python}
\label{subsec:rust-runner}

Rust und Python sind nicht über eine Sprachbindung wie PyO3 gekoppelt, sondern
lose über einen Subprozess-Aufruf und eine JSON-Datei verbunden. Die Funktion
\texttt{run\_rust\_and\_load\_json} kapselt diese Brücke. Sie prüft zunächst, ob
die Binary existiert, und bestimmt das Arbeitsverzeichnis anhand der
\texttt{input\_config.json}. Eine alte Ausgabedatei wird vorher gelöscht, damit
kein veraltetes Ergebnis eingelesen wird. Nach dem Lauf wertet die Funktion den
Rückgabecode aus; schlägt Rust fehl, landen \texttt{stdout} und \texttt{stderr} in
der Fehlermeldung und erleichtern die Diagnose. Der Vorteil der losen Kopplung
liegt im verteilten Betrieb: Beide Dienste lassen sich getrennt kompilieren,
testen und auf mehrere Cluster-Knoten verteilen.

\subsection{Signalverarbeitung mit SciPy}
\label{subsec:scipy}

Das Modul \texttt{convolution.py} bildet den Kern des Python-Dienstes. Es wandelt
die rohen Verzögerungs- und Druckdaten in eine abtastgenaue Impulsantwort um und
faltet damit das Trockensignal.

\subsubsection{Diskretisierung der Impulsantwort}
\label{subsubsec:build-ir}

Die Rust-Engine liefert eine endliche Menge diskreter Strahl-Treffer, jeweils
bestehend aus einer Ankunftsverzögerung $\tau_j$ (in Sekunden) und einem skalaren
Restdruck $p_j$. Die Funktion \texttt{build\_impulse\_response} quantisiert jeden
Treffer auf den nächstgelegenen Abtastindex $i_j = \operatorname{round}(\tau_j
\cdot f_s)$ und akkumuliert gleichzeitige Treffer zur diskreten Impulsantwort:
\[
    h[i] \;=\; \sum_{j} p_j \cdot \bigl[\,i_j = i\,\bigr].
\]
Die Akkumulation erfolgt mit \texttt{numpy.add.at}. Anders als die einfache
indizierte Zuweisung, die bei doppelten Indizes nur den letzten Wert behält,
summiert \texttt{add.at} gleichzeitige Treffer korrekt auf. Vor der Berechnung
prüft die Funktion die Eingaben an der Systemgrenze: übereinstimmende
Array-Formen, Eindimensionalität, nicht-negative Verzögerungen und eine positive
Abtastrate.

\begin{lstlisting}[style=cemdsp, caption={Diskretisierung der Ray-Hits zur Impulsantwort.}]
indices = np.round(delays * sample_rate).astype(np.int64)
ir_length = int(indices.max()) + 1
ir = np.zeros(ir_length, dtype=np.float64)
np.add.at(ir, indices, pressures)   # akkumuliert gleichzeitige Treffer
return ir
\end{lstlisting}

\subsubsection{Faltung mit \texttt{oaconvolve}}
\label{subsubsec:oaconvolve}

Die eigentliche Faltung nutzt \texttt{scipy.signal.oaconvolve}, die
Overlap-Add-basierte FFT-Faltung aus Abschnitt~2.3. Der Modus \texttt{"full"}
liefert das vollständige Ergebnis der Länge $N + M - 1$ und bewahrt damit den
kompletten Nachhall. Mehrkanaliges (etwa Stereo-)Audio wird kanalweise gefaltet
und danach wieder zusammengesetzt:

\begin{lstlisting}[style=cemdsp, caption={Mono- und Mehrkanal-Faltung.}]
if dry_audio.ndim == 1:
    wet = signal.oaconvolve(dry_audio, ir, mode="full")
elif dry_audio.ndim == 2:
    channels = [
        signal.oaconvolve(dry_audio[:, c], ir, mode="full")
        for c in range(dry_audio.shape[1])
    ]
    wet = np.stack(channels, axis=1)
return wet.astype(np.float32)
\end{lstlisting}

Vor der Faltung prüft \texttt{convolve} das Schema von \texttt{ir\_data} und
vergleicht die Abtastraten von Audio und Impulsantwort. Das Ergebnis gibt die
Funktion roh und unnormalisiert zurück. Die Lautstärke regelt erst der Export
(Abschnitt~\ref{subsec:rendering}); so bleibt die Faltung frei von Pegelfragen.

\subsubsection{Frequenzabhängige Materialdämpfung}
\label{subsubsec:material}

Reale Materialien dämpfen Schall frequenzabhängig. Ein Teppich schluckt hohe
Frequenzen viel stärker als tiefe, Beton reflektiert dagegen fast breitbandig.
Die Rust-Engine reduziert die Absorption aber schon auf einen einzigen skalaren
Restdruck $p_j$ pro Treffer, sodass die spektrale Information an dieser
Systemgrenze verloren geht. Damit der Nachhall trotzdem frequenzabhängig wird,
ergänzt \texttt{convolution.py} ein Absorptionsmodell im Frequenzbereich.

Grundlage sind die Oktavband-Absorptionskoeffizienten $\alpha_b$ aus der Datei
\texttt{materials.json}, die auch die Rust-Engine speist. Physik und
Signalverarbeitung greifen damit auf dieselbe Datenbasis zu (single source of
truth). Jedes Material ist über sechs Oktavbänder mit den Mittenfrequenzen
$125,\,250,\,500,\,1000,\,2000$ und $4000$~Hz beschrieben; ein hoher Koeffizient
bedeutet starke Dämpfung im jeweiligen Band.

Das Modell beruht auf einer einfachen Überlegung: Ein Tap, der zum Zeitpunkt $t$
eintrifft, hat näherungsweise $n(t) = t / \tau_{\mathrm{mfp}}$ Reflexionen hinter
sich, wobei $\tau_{\mathrm{mfp}}$ die mittlere Zeit zwischen zwei Reflexionen
angibt (die mittlere freie Weglänge in Zeiteinheiten). Nach $n$ Reflexionen
verbleibt in jedem Oktavband $b$ der Energieanteil
\[
    g_b(t) \;=\; \bigl(1 - \alpha_b\bigr)^{\,n(t)} \;=\; \bigl(1 - \alpha_b\bigr)^{\,t / \tau_{\mathrm{mfp}}}.
\]
Bänder mit hohem $\alpha_b$, etwa die Höhen bei Teppich, klingen dadurch deutlich
schneller ab als solche mit niedrigem $\alpha_b$. Der Nachhall wird über die Zeit
zunehmend dunkler.

Im Code setzt eine Oktavband-Filterbank dieses Modell um. Die Impulsantwort wird
über Butterworth-Bandpässe (nullphasig via \texttt{sosfiltfilt}) in ihre sechs
Bänder zerlegt, jedes Band mit der zeitabhängigen Hüllkurve $g_b(t)$ gewichtet und
die Bänder anschließend wieder aufsummiert.

\begin{lstlisting}[style=cemdsp, caption={Frequenzabhaengige Material-Daempfung der Impulsantwort.}]
centers, materials = _load_materials(str(_MATERIALS_JSON))
nyquist = sample_rate / 2.0
reflections = (np.arange(ir.size) / sample_rate) / mean_free_path_s

shaped = np.zeros_like(ir)
for (low, high), alpha in zip(_octave_band_edges(centers), materials[material]):
    band = _bandpass(ir, low, high, nyquist)      # ein Oktavband isolieren
    gain = np.power(max(1e-6, 1.0 - alpha), reflections)
    shaped += band * gain                          # zeitabhaengig daempfen
return shaped
\end{lstlisting}

\paragraph{Verifikation.} Ein reproduzierbarer Selbsttest prüft das Modell. Er
dämpft dieselbe dichte Impulsantwort einmal mit \texttt{carpet} und einmal mit
\texttt{concrete} und misst jeweils den Energieanteil oberhalb von $2$~kHz. Bei
Teppich bleiben davon nur $0{,}39$ übrig, bei Beton $0{,}71$; der Teppich schluckt
die Höhen also deutlich stärker. Kippt diese Reihenfolge, schlägt der Test fehl.

\paragraph{Grenzen des Modells.} Das Verfahren modelliert ein globales
Raummaterial, nicht die tatsächliche Materialfolge eines einzelnen Strahlpfads.
Eine Trennung nach Oberfläche, also die Aussage „dieser Tap stammt von Teppich,
jener von Beton", ist im Python-Dienst nicht möglich, solange die Rust-Engine nur
einen skalaren Druck pro Treffer liefert. Dafür müsste die Engine schon bandweise
Drücke ($p_{j,b}$ statt $p_j$) exportieren. Kapitel~7 führt das als mögliche
Erweiterung.

\subsection{Audio-Rendering und Export}
\label{subsec:rendering}

Das Ein- und Auslesen der Audiodateien sowie die Normalisierung gegen Clipping
übernimmt \texttt{wav\_handler.py}; die Details stehen in der Teildokumentation zum
Audio-I/O. Für den DSP-Teil zählt nur der Vertrag zwischen den Stufen.
\texttt{convolve} liefert ein rohes, womöglich übersteuertes Signal;
\texttt{write\_wav} misst dessen maximale Amplitude und normalisiert erst, wenn
sie den Bereich $[-1{,}0,\,1{,}0]$ verlässt. So bleibt Clipping aus, und die
Faltung muss sich nie um Pegel kümmern.

%   3. Der Master braucht nur diese Pakete (Dokumentation.tex hat beide schon):
%        \usepackage{amsmath}
%        \usepackage{listings}
%      Kein xcolor/babel noetig: Umlaute laufen unter UTF-8, deutsche
%      Anfuehrungszeichen sind literal „..." gesetzt.
%
% STANDALONE-VORSCHAU (zum alleine Kompilieren einfach drumherum setzen):
%   \documentclass{article}\usepackage{amsmath}\usepackage{listings}
%   \begin{document}
%   % =====================================================================
% Kapitel 5 - Digitale Signalverarbeitung und Orchestrierung (Python)
% Autor: Cem Acar
%
% EINBINDEN ins Hauptdokument:
%   1. Dieses File nach latexDocs/ legen.
%   2. Im Master (Dokumentation.tex) an der richtigen Stelle einfuegen:
%        \input{Kapitel5_DSP_Cem}
%   3. Der Master braucht nur diese Pakete (Dokumentation.tex hat beide schon):
%        \usepackage{amsmath}
%        \usepackage{listings}
%      Kein xcolor/babel noetig: Umlaute laufen unter UTF-8, deutsche
%      Anfuehrungszeichen sind literal „..." gesetzt.
%
% STANDALONE-VORSCHAU (zum alleine Kompilieren einfach drumherum setzen):
%   \documentclass{article}\usepackage{amsmath}\usepackage{listings}
%   \begin{document}
%   \input{Kapitel5_DSP_Cem}
%   \end{document}
% =====================================================================

% Listing-Stil fuer dieses Kapitel (nur listings-Kern, keine Farben noetig)
\lstdefinestyle{cemdsp}{
    language=Python,
    basicstyle=\ttfamily\small,
    showstringspaces=false,
    breaklines=true,
    numbers=left,
    numberstyle=\tiny,
    frame=single,
    tabsize=4
}

%======================================================================
\section{Digitale Signalverarbeitung und Orchestrierung (Python)}
\label{sec:python-dsp}

Die Rust-Engine berechnet die Physik der Schallausbreitung, der Python-Dienst
übernimmt die Signalverarbeitung und die Orchestrierung der Pipeline. Warum die
Physik in Rust und die Signalverarbeitung in Python mit NumPy und SciPy liegt,
ist in Kapitel~3 begründet. Die signaltheoretischen Grundlagen, also
Impulsantwort, Faltungsintegral und schnelle Faltung über die FFT, behandelt
Abschnitt~2.3. Dieses Kapitel zeigt, wie diese Grundlagen im Code umgesetzt sind,
und folgt dabei dem Datenfluss durch die Pipeline.

\subsection{Pipeline-Management durch \texttt{main.py}}
\label{subsec:pipeline}

\texttt{main.py} führt die vier Verarbeitungsschritte nacheinander aus. Tritt in
einem Schritt ein Fehler auf, bricht die Pipeline sofort mit einer Meldung ab,
statt mit ungültigen Zwischendaten weiterzurechnen (fail-fast).

\begin{enumerate}
    \item \textbf{Physik-Engine starten:} \texttt{run\_rust\_and\_load\_json} ruft
    die vorkompilierte Rust-Binary als Subprozess auf und lädt deren Ausgabe
    \texttt{ir\_output.json} als Python-Dictionary (\texttt{ir\_data}).
    \item \textbf{Trockenaudio laden:} \texttt{read\_wav} liest die
    Eingangsdatei und liefert Audiodaten samt Abtastrate.
    \item \textbf{Faltung:} \texttt{convolve} wendet die aus \texttt{ir\_data}
    erzeugte Impulsantwort auf das Trockensignal an.
    \item \textbf{Export:} \texttt{write\_wav} normalisiert das Ergebnis bei
    Bedarf und schreibt die fertige Ausgabedatei.
\end{enumerate}

Der Binary-Name wird plattformabhängig gewählt: unter Windows mit der Endung
\texttt{.exe}, unter macOS und Linux ohne. Stimmt die Abtastrate des Trockenaudios
nicht mit der Rate der Engine überein, gibt \texttt{main.py} eine Warnung aus,
weil der Hall sonst verstimmt klingen kann.

\subsection{Die Prozessbrücke zwischen Rust und Python}
\label{subsec:rust-runner}

Rust und Python sind nicht über eine Sprachbindung wie PyO3 gekoppelt, sondern
lose über einen Subprozess-Aufruf und eine JSON-Datei verbunden. Die Funktion
\texttt{run\_rust\_and\_load\_json} kapselt diese Brücke. Sie prüft zunächst, ob
die Binary existiert, und bestimmt das Arbeitsverzeichnis anhand der
\texttt{input\_config.json}. Eine alte Ausgabedatei wird vorher gelöscht, damit
kein veraltetes Ergebnis eingelesen wird. Nach dem Lauf wertet die Funktion den
Rückgabecode aus; schlägt Rust fehl, landen \texttt{stdout} und \texttt{stderr} in
der Fehlermeldung und erleichtern die Diagnose. Der Vorteil der losen Kopplung
liegt im verteilten Betrieb: Beide Dienste lassen sich getrennt kompilieren,
testen und auf mehrere Cluster-Knoten verteilen.

\subsection{Signalverarbeitung mit SciPy}
\label{subsec:scipy}

Das Modul \texttt{convolution.py} bildet den Kern des Python-Dienstes. Es wandelt
die rohen Verzögerungs- und Druckdaten in eine abtastgenaue Impulsantwort um und
faltet damit das Trockensignal.

\subsubsection{Diskretisierung der Impulsantwort}
\label{subsubsec:build-ir}

Die Rust-Engine liefert eine endliche Menge diskreter Strahl-Treffer, jeweils
bestehend aus einer Ankunftsverzögerung $\tau_j$ (in Sekunden) und einem skalaren
Restdruck $p_j$. Die Funktion \texttt{build\_impulse\_response} quantisiert jeden
Treffer auf den nächstgelegenen Abtastindex $i_j = \operatorname{round}(\tau_j
\cdot f_s)$ und akkumuliert gleichzeitige Treffer zur diskreten Impulsantwort:
\[
    h[i] \;=\; \sum_{j} p_j \cdot \bigl[\,i_j = i\,\bigr].
\]
Die Akkumulation erfolgt mit \texttt{numpy.add.at}. Anders als die einfache
indizierte Zuweisung, die bei doppelten Indizes nur den letzten Wert behält,
summiert \texttt{add.at} gleichzeitige Treffer korrekt auf. Vor der Berechnung
prüft die Funktion die Eingaben an der Systemgrenze: übereinstimmende
Array-Formen, Eindimensionalität, nicht-negative Verzögerungen und eine positive
Abtastrate.

\begin{lstlisting}[style=cemdsp, caption={Diskretisierung der Ray-Hits zur Impulsantwort.}]
indices = np.round(delays * sample_rate).astype(np.int64)
ir_length = int(indices.max()) + 1
ir = np.zeros(ir_length, dtype=np.float64)
np.add.at(ir, indices, pressures)   # akkumuliert gleichzeitige Treffer
return ir
\end{lstlisting}

\subsubsection{Faltung mit \texttt{oaconvolve}}
\label{subsubsec:oaconvolve}

Die eigentliche Faltung nutzt \texttt{scipy.signal.oaconvolve}, die
Overlap-Add-basierte FFT-Faltung aus Abschnitt~2.3. Der Modus \texttt{"full"}
liefert das vollständige Ergebnis der Länge $N + M - 1$ und bewahrt damit den
kompletten Nachhall. Mehrkanaliges (etwa Stereo-)Audio wird kanalweise gefaltet
und danach wieder zusammengesetzt:

\begin{lstlisting}[style=cemdsp, caption={Mono- und Mehrkanal-Faltung.}]
if dry_audio.ndim == 1:
    wet = signal.oaconvolve(dry_audio, ir, mode="full")
elif dry_audio.ndim == 2:
    channels = [
        signal.oaconvolve(dry_audio[:, c], ir, mode="full")
        for c in range(dry_audio.shape[1])
    ]
    wet = np.stack(channels, axis=1)
return wet.astype(np.float32)
\end{lstlisting}

Vor der Faltung prüft \texttt{convolve} das Schema von \texttt{ir\_data} und
vergleicht die Abtastraten von Audio und Impulsantwort. Das Ergebnis gibt die
Funktion roh und unnormalisiert zurück. Die Lautstärke regelt erst der Export
(Abschnitt~\ref{subsec:rendering}); so bleibt die Faltung frei von Pegelfragen.

\subsubsection{Frequenzabhängige Materialdämpfung}
\label{subsubsec:material}

Reale Materialien dämpfen Schall frequenzabhängig. Ein Teppich schluckt hohe
Frequenzen viel stärker als tiefe, Beton reflektiert dagegen fast breitbandig.
Die Rust-Engine reduziert die Absorption aber schon auf einen einzigen skalaren
Restdruck $p_j$ pro Treffer, sodass die spektrale Information an dieser
Systemgrenze verloren geht. Damit der Nachhall trotzdem frequenzabhängig wird,
ergänzt \texttt{convolution.py} ein Absorptionsmodell im Frequenzbereich.

Grundlage sind die Oktavband-Absorptionskoeffizienten $\alpha_b$ aus der Datei
\texttt{materials.json}, die auch die Rust-Engine speist. Physik und
Signalverarbeitung greifen damit auf dieselbe Datenbasis zu (single source of
truth). Jedes Material ist über sechs Oktavbänder mit den Mittenfrequenzen
$125,\,250,\,500,\,1000,\,2000$ und $4000$~Hz beschrieben; ein hoher Koeffizient
bedeutet starke Dämpfung im jeweiligen Band.

Das Modell beruht auf einer einfachen Überlegung: Ein Tap, der zum Zeitpunkt $t$
eintrifft, hat näherungsweise $n(t) = t / \tau_{\mathrm{mfp}}$ Reflexionen hinter
sich, wobei $\tau_{\mathrm{mfp}}$ die mittlere Zeit zwischen zwei Reflexionen
angibt (die mittlere freie Weglänge in Zeiteinheiten). Nach $n$ Reflexionen
verbleibt in jedem Oktavband $b$ der Energieanteil
\[
    g_b(t) \;=\; \bigl(1 - \alpha_b\bigr)^{\,n(t)} \;=\; \bigl(1 - \alpha_b\bigr)^{\,t / \tau_{\mathrm{mfp}}}.
\]
Bänder mit hohem $\alpha_b$, etwa die Höhen bei Teppich, klingen dadurch deutlich
schneller ab als solche mit niedrigem $\alpha_b$. Der Nachhall wird über die Zeit
zunehmend dunkler.

Im Code setzt eine Oktavband-Filterbank dieses Modell um. Die Impulsantwort wird
über Butterworth-Bandpässe (nullphasig via \texttt{sosfiltfilt}) in ihre sechs
Bänder zerlegt, jedes Band mit der zeitabhängigen Hüllkurve $g_b(t)$ gewichtet und
die Bänder anschließend wieder aufsummiert.

\begin{lstlisting}[style=cemdsp, caption={Frequenzabhaengige Material-Daempfung der Impulsantwort.}]
centers, materials = _load_materials(str(_MATERIALS_JSON))
nyquist = sample_rate / 2.0
reflections = (np.arange(ir.size) / sample_rate) / mean_free_path_s

shaped = np.zeros_like(ir)
for (low, high), alpha in zip(_octave_band_edges(centers), materials[material]):
    band = _bandpass(ir, low, high, nyquist)      # ein Oktavband isolieren
    gain = np.power(max(1e-6, 1.0 - alpha), reflections)
    shaped += band * gain                          # zeitabhaengig daempfen
return shaped
\end{lstlisting}

\paragraph{Verifikation.} Ein reproduzierbarer Selbsttest prüft das Modell. Er
dämpft dieselbe dichte Impulsantwort einmal mit \texttt{carpet} und einmal mit
\texttt{concrete} und misst jeweils den Energieanteil oberhalb von $2$~kHz. Bei
Teppich bleiben davon nur $0{,}39$ übrig, bei Beton $0{,}71$; der Teppich schluckt
die Höhen also deutlich stärker. Kippt diese Reihenfolge, schlägt der Test fehl.

\paragraph{Grenzen des Modells.} Das Verfahren modelliert ein globales
Raummaterial, nicht die tatsächliche Materialfolge eines einzelnen Strahlpfads.
Eine Trennung nach Oberfläche, also die Aussage „dieser Tap stammt von Teppich,
jener von Beton", ist im Python-Dienst nicht möglich, solange die Rust-Engine nur
einen skalaren Druck pro Treffer liefert. Dafür müsste die Engine schon bandweise
Drücke ($p_{j,b}$ statt $p_j$) exportieren. Kapitel~7 führt das als mögliche
Erweiterung.

\subsection{Audio-Rendering und Export}
\label{subsec:rendering}

Das Ein- und Auslesen der Audiodateien sowie die Normalisierung gegen Clipping
übernimmt \texttt{wav\_handler.py}; die Details stehen in der Teildokumentation zum
Audio-I/O. Für den DSP-Teil zählt nur der Vertrag zwischen den Stufen.
\texttt{convolve} liefert ein rohes, womöglich übersteuertes Signal;
\texttt{write\_wav} misst dessen maximale Amplitude und normalisiert erst, wenn
sie den Bereich $[-1{,}0,\,1{,}0]$ verlässt. So bleibt Clipping aus, und die
Faltung muss sich nie um Pegel kümmern.

%   \end{document}
% =====================================================================

% Listing-Stil fuer dieses Kapitel (nur listings-Kern, keine Farben noetig)
\lstdefinestyle{cemdsp}{
    language=Python,
    basicstyle=\ttfamily\small,
    showstringspaces=false,
    breaklines=true,
    numbers=left,
    numberstyle=\tiny,
    frame=single,
    tabsize=4
}

%======================================================================
\section{Digitale Signalverarbeitung und Orchestrierung (Python)}
\label{sec:python-dsp}

Die Rust-Engine berechnet die Physik der Schallausbreitung, der Python-Dienst
übernimmt die Signalverarbeitung und die Orchestrierung der Pipeline. Warum die
Physik in Rust und die Signalverarbeitung in Python mit NumPy und SciPy liegt,
ist in Kapitel~3 begründet. Die signaltheoretischen Grundlagen, also
Impulsantwort, Faltungsintegral und schnelle Faltung über die FFT, behandelt
Abschnitt~2.3. Dieses Kapitel zeigt, wie diese Grundlagen im Code umgesetzt sind,
und folgt dabei dem Datenfluss durch die Pipeline.

\subsection{Pipeline-Management durch \texttt{main.py}}
\label{subsec:pipeline}

\texttt{main.py} führt die vier Verarbeitungsschritte nacheinander aus. Tritt in
einem Schritt ein Fehler auf, bricht die Pipeline sofort mit einer Meldung ab,
statt mit ungültigen Zwischendaten weiterzurechnen (fail-fast).

\begin{enumerate}
    \item \textbf{Physik-Engine starten:} \texttt{run\_rust\_and\_load\_json} ruft
    die vorkompilierte Rust-Binary als Subprozess auf und lädt deren Ausgabe
    \texttt{ir\_output.json} als Python-Dictionary (\texttt{ir\_data}).
    \item \textbf{Trockenaudio laden:} \texttt{read\_wav} liest die
    Eingangsdatei und liefert Audiodaten samt Abtastrate.
    \item \textbf{Faltung:} \texttt{convolve} wendet die aus \texttt{ir\_data}
    erzeugte Impulsantwort auf das Trockensignal an.
    \item \textbf{Export:} \texttt{write\_wav} normalisiert das Ergebnis bei
    Bedarf und schreibt die fertige Ausgabedatei.
\end{enumerate}

Der Binary-Name wird plattformabhängig gewählt: unter Windows mit der Endung
\texttt{.exe}, unter macOS und Linux ohne. Stimmt die Abtastrate des Trockenaudios
nicht mit der Rate der Engine überein, gibt \texttt{main.py} eine Warnung aus,
weil der Hall sonst verstimmt klingen kann.

\subsection{Die Prozessbrücke zwischen Rust und Python}
\label{subsec:rust-runner}

Rust und Python sind nicht über eine Sprachbindung wie PyO3 gekoppelt, sondern
lose über einen Subprozess-Aufruf und eine JSON-Datei verbunden. Die Funktion
\texttt{run\_rust\_and\_load\_json} kapselt diese Brücke. Sie prüft zunächst, ob
die Binary existiert, und bestimmt das Arbeitsverzeichnis anhand der
\texttt{input\_config.json}. Eine alte Ausgabedatei wird vorher gelöscht, damit
kein veraltetes Ergebnis eingelesen wird. Nach dem Lauf wertet die Funktion den
Rückgabecode aus; schlägt Rust fehl, landen \texttt{stdout} und \texttt{stderr} in
der Fehlermeldung und erleichtern die Diagnose. Der Vorteil der losen Kopplung
liegt im verteilten Betrieb: Beide Dienste lassen sich getrennt kompilieren,
testen und auf mehrere Cluster-Knoten verteilen.

\subsection{Signalverarbeitung mit SciPy}
\label{subsec:scipy}

Das Modul \texttt{convolution.py} bildet den Kern des Python-Dienstes. Es wandelt
die rohen Verzögerungs- und Druckdaten in eine abtastgenaue Impulsantwort um und
faltet damit das Trockensignal.

\subsubsection{Diskretisierung der Impulsantwort}
\label{subsubsec:build-ir}

Die Rust-Engine liefert eine endliche Menge diskreter Strahl-Treffer, jeweils
bestehend aus einer Ankunftsverzögerung $\tau_j$ (in Sekunden) und einem skalaren
Restdruck $p_j$. Die Funktion \texttt{build\_impulse\_response} quantisiert jeden
Treffer auf den nächstgelegenen Abtastindex $i_j = \operatorname{round}(\tau_j
\cdot f_s)$ und akkumuliert gleichzeitige Treffer zur diskreten Impulsantwort:
\[
    h[i] \;=\; \sum_{j} p_j \cdot \bigl[\,i_j = i\,\bigr].
\]
Die Akkumulation erfolgt mit \texttt{numpy.add.at}. Anders als die einfache
indizierte Zuweisung, die bei doppelten Indizes nur den letzten Wert behält,
summiert \texttt{add.at} gleichzeitige Treffer korrekt auf. Vor der Berechnung
prüft die Funktion die Eingaben an der Systemgrenze: übereinstimmende
Array-Formen, Eindimensionalität, nicht-negative Verzögerungen und eine positive
Abtastrate.

\begin{lstlisting}[style=cemdsp, caption={Diskretisierung der Ray-Hits zur Impulsantwort.}]
indices = np.round(delays * sample_rate).astype(np.int64)
ir_length = int(indices.max()) + 1
ir = np.zeros(ir_length, dtype=np.float64)
np.add.at(ir, indices, pressures)   # akkumuliert gleichzeitige Treffer
return ir
\end{lstlisting}

\subsubsection{Faltung mit \texttt{oaconvolve}}
\label{subsubsec:oaconvolve}

Die eigentliche Faltung nutzt \texttt{scipy.signal.oaconvolve}, die
Overlap-Add-basierte FFT-Faltung aus Abschnitt~2.3. Der Modus \texttt{"full"}
liefert das vollständige Ergebnis der Länge $N + M - 1$ und bewahrt damit den
kompletten Nachhall. Mehrkanaliges (etwa Stereo-)Audio wird kanalweise gefaltet
und danach wieder zusammengesetzt:

\begin{lstlisting}[style=cemdsp, caption={Mono- und Mehrkanal-Faltung.}]
if dry_audio.ndim == 1:
    wet = signal.oaconvolve(dry_audio, ir, mode="full")
elif dry_audio.ndim == 2:
    channels = [
        signal.oaconvolve(dry_audio[:, c], ir, mode="full")
        for c in range(dry_audio.shape[1])
    ]
    wet = np.stack(channels, axis=1)
return wet.astype(np.float32)
\end{lstlisting}

Vor der Faltung prüft \texttt{convolve} das Schema von \texttt{ir\_data} und
vergleicht die Abtastraten von Audio und Impulsantwort. Das Ergebnis gibt die
Funktion roh und unnormalisiert zurück. Die Lautstärke regelt erst der Export
(Abschnitt~\ref{subsec:rendering}); so bleibt die Faltung frei von Pegelfragen.

\subsubsection{Frequenzabhängige Materialdämpfung}
\label{subsubsec:material}

Reale Materialien dämpfen Schall frequenzabhängig. Ein Teppich schluckt hohe
Frequenzen viel stärker als tiefe, Beton reflektiert dagegen fast breitbandig.
Die Rust-Engine reduziert die Absorption aber schon auf einen einzigen skalaren
Restdruck $p_j$ pro Treffer, sodass die spektrale Information an dieser
Systemgrenze verloren geht. Damit der Nachhall trotzdem frequenzabhängig wird,
ergänzt \texttt{convolution.py} ein Absorptionsmodell im Frequenzbereich.

Grundlage sind die Oktavband-Absorptionskoeffizienten $\alpha_b$ aus der Datei
\texttt{materials.json}, die auch die Rust-Engine speist. Physik und
Signalverarbeitung greifen damit auf dieselbe Datenbasis zu (single source of
truth). Jedes Material ist über sechs Oktavbänder mit den Mittenfrequenzen
$125,\,250,\,500,\,1000,\,2000$ und $4000$~Hz beschrieben; ein hoher Koeffizient
bedeutet starke Dämpfung im jeweiligen Band.

Das Modell beruht auf einer einfachen Überlegung: Ein Tap, der zum Zeitpunkt $t$
eintrifft, hat näherungsweise $n(t) = t / \tau_{\mathrm{mfp}}$ Reflexionen hinter
sich, wobei $\tau_{\mathrm{mfp}}$ die mittlere Zeit zwischen zwei Reflexionen
angibt (die mittlere freie Weglänge in Zeiteinheiten). Nach $n$ Reflexionen
verbleibt in jedem Oktavband $b$ der Energieanteil
\[
    g_b(t) \;=\; \bigl(1 - \alpha_b\bigr)^{\,n(t)} \;=\; \bigl(1 - \alpha_b\bigr)^{\,t / \tau_{\mathrm{mfp}}}.
\]
Bänder mit hohem $\alpha_b$, etwa die Höhen bei Teppich, klingen dadurch deutlich
schneller ab als solche mit niedrigem $\alpha_b$. Der Nachhall wird über die Zeit
zunehmend dunkler.

Im Code setzt eine Oktavband-Filterbank dieses Modell um. Die Impulsantwort wird
über Butterworth-Bandpässe (nullphasig via \texttt{sosfiltfilt}) in ihre sechs
Bänder zerlegt, jedes Band mit der zeitabhängigen Hüllkurve $g_b(t)$ gewichtet und
die Bänder anschließend wieder aufsummiert.

\begin{lstlisting}[style=cemdsp, caption={Frequenzabhaengige Material-Daempfung der Impulsantwort.}]
centers, materials = _load_materials(str(_MATERIALS_JSON))
nyquist = sample_rate / 2.0
reflections = (np.arange(ir.size) / sample_rate) / mean_free_path_s

shaped = np.zeros_like(ir)
for (low, high), alpha in zip(_octave_band_edges(centers), materials[material]):
    band = _bandpass(ir, low, high, nyquist)      # ein Oktavband isolieren
    gain = np.power(max(1e-6, 1.0 - alpha), reflections)
    shaped += band * gain                          # zeitabhaengig daempfen
return shaped
\end{lstlisting}

\paragraph{Verifikation.} Ein reproduzierbarer Selbsttest prüft das Modell. Er
dämpft dieselbe dichte Impulsantwort einmal mit \texttt{carpet} und einmal mit
\texttt{concrete} und misst jeweils den Energieanteil oberhalb von $2$~kHz. Bei
Teppich bleiben davon nur $0{,}39$ übrig, bei Beton $0{,}71$; der Teppich schluckt
die Höhen also deutlich stärker. Kippt diese Reihenfolge, schlägt der Test fehl.

\paragraph{Grenzen des Modells.} Das Verfahren modelliert ein globales
Raummaterial, nicht die tatsächliche Materialfolge eines einzelnen Strahlpfads.
Eine Trennung nach Oberfläche, also die Aussage „dieser Tap stammt von Teppich,
jener von Beton", ist im Python-Dienst nicht möglich, solange die Rust-Engine nur
einen skalaren Druck pro Treffer liefert. Dafür müsste die Engine schon bandweise
Drücke ($p_{j,b}$ statt $p_j$) exportieren. Kapitel~7 führt das als mögliche
Erweiterung.

\subsection{Audio-Rendering und Export}
\label{subsec:rendering}

Das Ein- und Auslesen der Audiodateien sowie die Normalisierung gegen Clipping
übernimmt \texttt{wav\_handler.py}; die Details stehen in der Teildokumentation zum
Audio-I/O. Für den DSP-Teil zählt nur der Vertrag zwischen den Stufen.
\texttt{convolve} liefert ein rohes, womöglich übersteuertes Signal;
\texttt{write\_wav} misst dessen maximale Amplitude und normalisiert erst, wenn
sie den Bereich $[-1{,}0,\,1{,}0]$ verlässt. So bleibt Clipping aus, und die
Faltung muss sich nie um Pegel kümmern.

%   \end{document}
% =====================================================================

% Listing-Stil fuer dieses Kapitel (nur listings-Kern, keine Farben noetig)
\lstdefinestyle{cemdsp}{
    language=Python,
    basicstyle=\ttfamily\small,
    showstringspaces=false,
    breaklines=true,
    numbers=left,
    numberstyle=\tiny,
    frame=single,
    tabsize=4
}

%======================================================================
\section{Digitale Signalverarbeitung und Orchestrierung (Python)}
\label{sec:python-dsp}

Die Rust-Engine berechnet die Physik der Schallausbreitung, der Python-Dienst
übernimmt die Signalverarbeitung und die Orchestrierung der Pipeline. Warum die
Physik in Rust und die Signalverarbeitung in Python mit NumPy und SciPy liegt,
ist in Kapitel~3 begründet. Die signaltheoretischen Grundlagen, also
Impulsantwort, Faltungsintegral und schnelle Faltung über die FFT, behandelt
Abschnitt~2.3. Dieses Kapitel zeigt, wie diese Grundlagen im Code umgesetzt sind,
und folgt dabei dem Datenfluss durch die Pipeline.

\subsection{Pipeline-Management durch \texttt{main.py}}
\label{subsec:pipeline}

\texttt{main.py} führt die vier Verarbeitungsschritte nacheinander aus. Tritt in
einem Schritt ein Fehler auf, bricht die Pipeline sofort mit einer Meldung ab,
statt mit ungültigen Zwischendaten weiterzurechnen (fail-fast).

\begin{enumerate}
    \item \textbf{Physik-Engine starten:} \texttt{run\_rust\_and\_load\_json} ruft
    die vorkompilierte Rust-Binary als Subprozess auf und lädt deren Ausgabe
    \texttt{ir\_output.json} als Python-Dictionary (\texttt{ir\_data}).
    \item \textbf{Trockenaudio laden:} \texttt{read\_wav} liest die
    Eingangsdatei und liefert Audiodaten samt Abtastrate.
    \item \textbf{Faltung:} \texttt{convolve} wendet die aus \texttt{ir\_data}
    erzeugte Impulsantwort auf das Trockensignal an.
    \item \textbf{Export:} \texttt{write\_wav} normalisiert das Ergebnis bei
    Bedarf und schreibt die fertige Ausgabedatei.
\end{enumerate}

Der Binary-Name wird plattformabhängig gewählt: unter Windows mit der Endung
\texttt{.exe}, unter macOS und Linux ohne. Stimmt die Abtastrate des Trockenaudios
nicht mit der Rate der Engine überein, gibt \texttt{main.py} eine Warnung aus,
weil der Hall sonst verstimmt klingen kann.

\subsection{Die Prozessbrücke zwischen Rust und Python}
\label{subsec:rust-runner}

Rust und Python sind nicht über eine Sprachbindung wie PyO3 gekoppelt, sondern
lose über einen Subprozess-Aufruf und eine JSON-Datei verbunden. Die Funktion
\texttt{run\_rust\_and\_load\_json} kapselt diese Brücke. Sie prüft zunächst, ob
die Binary existiert, und bestimmt das Arbeitsverzeichnis anhand der
\texttt{input\_config.json}. Eine alte Ausgabedatei wird vorher gelöscht, damit
kein veraltetes Ergebnis eingelesen wird. Nach dem Lauf wertet die Funktion den
Rückgabecode aus; schlägt Rust fehl, landen \texttt{stdout} und \texttt{stderr} in
der Fehlermeldung und erleichtern die Diagnose. Der Vorteil der losen Kopplung
liegt im verteilten Betrieb: Beide Dienste lassen sich getrennt kompilieren,
testen und auf mehrere Cluster-Knoten verteilen.

\subsection{Signalverarbeitung mit SciPy}
\label{subsec:scipy}

Das Modul \texttt{convolution.py} bildet den Kern des Python-Dienstes. Es wandelt
die rohen Verzögerungs- und Druckdaten in eine abtastgenaue Impulsantwort um und
faltet damit das Trockensignal.

\subsubsection{Diskretisierung der Impulsantwort}
\label{subsubsec:build-ir}

Die Rust-Engine liefert eine endliche Menge diskreter Strahl-Treffer, jeweils
bestehend aus einer Ankunftsverzögerung $\tau_j$ (in Sekunden) und einem skalaren
Restdruck $p_j$. Die Funktion \texttt{build\_impulse\_response} quantisiert jeden
Treffer auf den nächstgelegenen Abtastindex $i_j = \operatorname{round}(\tau_j
\cdot f_s)$ und akkumuliert gleichzeitige Treffer zur diskreten Impulsantwort:
\[
    h[i] \;=\; \sum_{j} p_j \cdot \bigl[\,i_j = i\,\bigr].
\]
Die Akkumulation erfolgt mit \texttt{numpy.add.at}. Anders als die einfache
indizierte Zuweisung, die bei doppelten Indizes nur den letzten Wert behält,
summiert \texttt{add.at} gleichzeitige Treffer korrekt auf. Vor der Berechnung
prüft die Funktion die Eingaben an der Systemgrenze: übereinstimmende
Array-Formen, Eindimensionalität, nicht-negative Verzögerungen und eine positive
Abtastrate.

\begin{lstlisting}[style=cemdsp, caption={Diskretisierung der Ray-Hits zur Impulsantwort.}]
indices = np.round(delays * sample_rate).astype(np.int64)
ir_length = int(indices.max()) + 1
ir = np.zeros(ir_length, dtype=np.float64)
np.add.at(ir, indices, pressures)   # akkumuliert gleichzeitige Treffer
return ir
\end{lstlisting}

\subsubsection{Faltung mit \texttt{oaconvolve}}
\label{subsubsec:oaconvolve}

Die eigentliche Faltung nutzt \texttt{scipy.signal.oaconvolve}, die
Overlap-Add-basierte FFT-Faltung aus Abschnitt~2.3. Der Modus \texttt{"full"}
liefert das vollständige Ergebnis der Länge $N + M - 1$ und bewahrt damit den
kompletten Nachhall. Mehrkanaliges (etwa Stereo-)Audio wird kanalweise gefaltet
und danach wieder zusammengesetzt:

\begin{lstlisting}[style=cemdsp, caption={Mono- und Mehrkanal-Faltung.}]
if dry_audio.ndim == 1:
    wet = signal.oaconvolve(dry_audio, ir, mode="full")
elif dry_audio.ndim == 2:
    channels = [
        signal.oaconvolve(dry_audio[:, c], ir, mode="full")
        for c in range(dry_audio.shape[1])
    ]
    wet = np.stack(channels, axis=1)
return wet.astype(np.float32)
\end{lstlisting}

Vor der Faltung prüft \texttt{convolve} das Schema von \texttt{ir\_data} und
vergleicht die Abtastraten von Audio und Impulsantwort. Das Ergebnis gibt die
Funktion roh und unnormalisiert zurück. Die Lautstärke regelt erst der Export
(Abschnitt~\ref{subsec:rendering}); so bleibt die Faltung frei von Pegelfragen.

\subsubsection{Frequenzabhängige Materialdämpfung}
\label{subsubsec:material}

Reale Materialien dämpfen Schall frequenzabhängig. Ein Teppich schluckt hohe
Frequenzen viel stärker als tiefe, Beton reflektiert dagegen fast breitbandig.
Die Rust-Engine reduziert die Absorption aber schon auf einen einzigen skalaren
Restdruck $p_j$ pro Treffer, sodass die spektrale Information an dieser
Systemgrenze verloren geht. Damit der Nachhall trotzdem frequenzabhängig wird,
ergänzt \texttt{convolution.py} ein Absorptionsmodell im Frequenzbereich.

Grundlage sind die Oktavband-Absorptionskoeffizienten $\alpha_b$ aus der Datei
\texttt{materials.json}, die auch die Rust-Engine speist. Physik und
Signalverarbeitung greifen damit auf dieselbe Datenbasis zu (single source of
truth). Jedes Material ist über sechs Oktavbänder mit den Mittenfrequenzen
$125,\,250,\,500,\,1000,\,2000$ und $4000$~Hz beschrieben; ein hoher Koeffizient
bedeutet starke Dämpfung im jeweiligen Band.

Das Modell beruht auf einer einfachen Überlegung: Ein Tap, der zum Zeitpunkt $t$
eintrifft, hat näherungsweise $n(t) = t / \tau_{\mathrm{mfp}}$ Reflexionen hinter
sich, wobei $\tau_{\mathrm{mfp}}$ die mittlere Zeit zwischen zwei Reflexionen
angibt (die mittlere freie Weglänge in Zeiteinheiten). Nach $n$ Reflexionen
verbleibt in jedem Oktavband $b$ der Energieanteil
\[
    g_b(t) \;=\; \bigl(1 - \alpha_b\bigr)^{\,n(t)} \;=\; \bigl(1 - \alpha_b\bigr)^{\,t / \tau_{\mathrm{mfp}}}.
\]
Bänder mit hohem $\alpha_b$, etwa die Höhen bei Teppich, klingen dadurch deutlich
schneller ab als solche mit niedrigem $\alpha_b$. Der Nachhall wird über die Zeit
zunehmend dunkler.

Im Code setzt eine Oktavband-Filterbank dieses Modell um. Die Impulsantwort wird
über Butterworth-Bandpässe (nullphasig via \texttt{sosfiltfilt}) in ihre sechs
Bänder zerlegt, jedes Band mit der zeitabhängigen Hüllkurve $g_b(t)$ gewichtet und
die Bänder anschließend wieder aufsummiert.

\begin{lstlisting}[style=cemdsp, caption={Frequenzabhaengige Material-Daempfung der Impulsantwort.}]
centers, materials = _load_materials(str(_MATERIALS_JSON))
nyquist = sample_rate / 2.0
reflections = (np.arange(ir.size) / sample_rate) / mean_free_path_s

shaped = np.zeros_like(ir)
for (low, high), alpha in zip(_octave_band_edges(centers), materials[material]):
    band = _bandpass(ir, low, high, nyquist)      # ein Oktavband isolieren
    gain = np.power(max(1e-6, 1.0 - alpha), reflections)
    shaped += band * gain                          # zeitabhaengig daempfen
return shaped
\end{lstlisting}

\paragraph{Verifikation.} Ein reproduzierbarer Selbsttest prüft das Modell. Er
dämpft dieselbe dichte Impulsantwort einmal mit \texttt{carpet} und einmal mit
\texttt{concrete} und misst jeweils den Energieanteil oberhalb von $2$~kHz. Bei
Teppich bleiben davon nur $0{,}39$ übrig, bei Beton $0{,}71$; der Teppich schluckt
die Höhen also deutlich stärker. Kippt diese Reihenfolge, schlägt der Test fehl.

\paragraph{Grenzen des Modells.} Das Verfahren modelliert ein globales
Raummaterial, nicht die tatsächliche Materialfolge eines einzelnen Strahlpfads.
Eine Trennung nach Oberfläche, also die Aussage „dieser Tap stammt von Teppich,
jener von Beton", ist im Python-Dienst nicht möglich, solange die Rust-Engine nur
einen skalaren Druck pro Treffer liefert. Dafür müsste die Engine schon bandweise
Drücke ($p_{j,b}$ statt $p_j$) exportieren. Kapitel~7 führt das als mögliche
Erweiterung.

\subsection{Audio-Rendering und Export}
\label{subsec:rendering}

Das Ein- und Auslesen der Audiodateien sowie die Normalisierung gegen Clipping
übernimmt \texttt{wav\_handler.py}; die Details stehen in der Teildokumentation zum
Audio-I/O. Für den DSP-Teil zählt nur der Vertrag zwischen den Stufen.
\texttt{convolve} liefert ein rohes, womöglich übersteuertes Signal;
\texttt{write\_wav} misst dessen maximale Amplitude und normalisiert erst, wenn
sie den Bereich $[-1{,}0,\,1{,}0]$ verlässt. So bleibt Clipping aus, und die
Faltung muss sich nie um Pegel kümmern.

%   3. Der Master braucht nur diese Pakete (Dokumentation.tex hat beide schon):
%        \usepackage{amsmath}
%        \usepackage{listings}
%      Kein xcolor/babel noetig: Umlaute laufen unter UTF-8, deutsche
%      Anfuehrungszeichen sind literal „..." gesetzt.
%
% STANDALONE-VORSCHAU (zum alleine Kompilieren einfach drumherum setzen):
%   \documentclass{article}\usepackage{amsmath}\usepackage{listings}
%   \begin{document}
%   % =====================================================================
% Kapitel 5 - Digitale Signalverarbeitung und Orchestrierung (Python)
% Autor: Cem Acar
%
% EINBINDEN ins Hauptdokument:
%   1. Dieses File nach latexDocs/ legen.
%   2. Im Master (Dokumentation.tex) an der richtigen Stelle einfuegen:
%        % =====================================================================
% Kapitel 5 - Digitale Signalverarbeitung und Orchestrierung (Python)
% Autor: Cem Acar
%
% EINBINDEN ins Hauptdokument:
%   1. Dieses File nach latexDocs/ legen.
%   2. Im Master (Dokumentation.tex) an der richtigen Stelle einfuegen:
%        % =====================================================================
% Kapitel 5 - Digitale Signalverarbeitung und Orchestrierung (Python)
% Autor: Cem Acar
%
% EINBINDEN ins Hauptdokument:
%   1. Dieses File nach latexDocs/ legen.
%   2. Im Master (Dokumentation.tex) an der richtigen Stelle einfuegen:
%        \input{Kapitel5_DSP_Cem}
%   3. Der Master braucht nur diese Pakete (Dokumentation.tex hat beide schon):
%        \usepackage{amsmath}
%        \usepackage{listings}
%      Kein xcolor/babel noetig: Umlaute laufen unter UTF-8, deutsche
%      Anfuehrungszeichen sind literal „..." gesetzt.
%
% STANDALONE-VORSCHAU (zum alleine Kompilieren einfach drumherum setzen):
%   \documentclass{article}\usepackage{amsmath}\usepackage{listings}
%   \begin{document}
%   \input{Kapitel5_DSP_Cem}
%   \end{document}
% =====================================================================

% Listing-Stil fuer dieses Kapitel (nur listings-Kern, keine Farben noetig)
\lstdefinestyle{cemdsp}{
    language=Python,
    basicstyle=\ttfamily\small,
    showstringspaces=false,
    breaklines=true,
    numbers=left,
    numberstyle=\tiny,
    frame=single,
    tabsize=4
}

%======================================================================
\section{Digitale Signalverarbeitung und Orchestrierung (Python)}
\label{sec:python-dsp}

Die Rust-Engine berechnet die Physik der Schallausbreitung, der Python-Dienst
übernimmt die Signalverarbeitung und die Orchestrierung der Pipeline. Warum die
Physik in Rust und die Signalverarbeitung in Python mit NumPy und SciPy liegt,
ist in Kapitel~3 begründet. Die signaltheoretischen Grundlagen, also
Impulsantwort, Faltungsintegral und schnelle Faltung über die FFT, behandelt
Abschnitt~2.3. Dieses Kapitel zeigt, wie diese Grundlagen im Code umgesetzt sind,
und folgt dabei dem Datenfluss durch die Pipeline.

\subsection{Pipeline-Management durch \texttt{main.py}}
\label{subsec:pipeline}

\texttt{main.py} führt die vier Verarbeitungsschritte nacheinander aus. Tritt in
einem Schritt ein Fehler auf, bricht die Pipeline sofort mit einer Meldung ab,
statt mit ungültigen Zwischendaten weiterzurechnen (fail-fast).

\begin{enumerate}
    \item \textbf{Physik-Engine starten:} \texttt{run\_rust\_and\_load\_json} ruft
    die vorkompilierte Rust-Binary als Subprozess auf und lädt deren Ausgabe
    \texttt{ir\_output.json} als Python-Dictionary (\texttt{ir\_data}).
    \item \textbf{Trockenaudio laden:} \texttt{read\_wav} liest die
    Eingangsdatei und liefert Audiodaten samt Abtastrate.
    \item \textbf{Faltung:} \texttt{convolve} wendet die aus \texttt{ir\_data}
    erzeugte Impulsantwort auf das Trockensignal an.
    \item \textbf{Export:} \texttt{write\_wav} normalisiert das Ergebnis bei
    Bedarf und schreibt die fertige Ausgabedatei.
\end{enumerate}

Der Binary-Name wird plattformabhängig gewählt: unter Windows mit der Endung
\texttt{.exe}, unter macOS und Linux ohne. Stimmt die Abtastrate des Trockenaudios
nicht mit der Rate der Engine überein, gibt \texttt{main.py} eine Warnung aus,
weil der Hall sonst verstimmt klingen kann.

\subsection{Die Prozessbrücke zwischen Rust und Python}
\label{subsec:rust-runner}

Rust und Python sind nicht über eine Sprachbindung wie PyO3 gekoppelt, sondern
lose über einen Subprozess-Aufruf und eine JSON-Datei verbunden. Die Funktion
\texttt{run\_rust\_and\_load\_json} kapselt diese Brücke. Sie prüft zunächst, ob
die Binary existiert, und bestimmt das Arbeitsverzeichnis anhand der
\texttt{input\_config.json}. Eine alte Ausgabedatei wird vorher gelöscht, damit
kein veraltetes Ergebnis eingelesen wird. Nach dem Lauf wertet die Funktion den
Rückgabecode aus; schlägt Rust fehl, landen \texttt{stdout} und \texttt{stderr} in
der Fehlermeldung und erleichtern die Diagnose. Der Vorteil der losen Kopplung
liegt im verteilten Betrieb: Beide Dienste lassen sich getrennt kompilieren,
testen und auf mehrere Cluster-Knoten verteilen.

\subsection{Signalverarbeitung mit SciPy}
\label{subsec:scipy}

Das Modul \texttt{convolution.py} bildet den Kern des Python-Dienstes. Es wandelt
die rohen Verzögerungs- und Druckdaten in eine abtastgenaue Impulsantwort um und
faltet damit das Trockensignal.

\subsubsection{Diskretisierung der Impulsantwort}
\label{subsubsec:build-ir}

Die Rust-Engine liefert eine endliche Menge diskreter Strahl-Treffer, jeweils
bestehend aus einer Ankunftsverzögerung $\tau_j$ (in Sekunden) und einem skalaren
Restdruck $p_j$. Die Funktion \texttt{build\_impulse\_response} quantisiert jeden
Treffer auf den nächstgelegenen Abtastindex $i_j = \operatorname{round}(\tau_j
\cdot f_s)$ und akkumuliert gleichzeitige Treffer zur diskreten Impulsantwort:
\[
    h[i] \;=\; \sum_{j} p_j \cdot \bigl[\,i_j = i\,\bigr].
\]
Die Akkumulation erfolgt mit \texttt{numpy.add.at}. Anders als die einfache
indizierte Zuweisung, die bei doppelten Indizes nur den letzten Wert behält,
summiert \texttt{add.at} gleichzeitige Treffer korrekt auf. Vor der Berechnung
prüft die Funktion die Eingaben an der Systemgrenze: übereinstimmende
Array-Formen, Eindimensionalität, nicht-negative Verzögerungen und eine positive
Abtastrate.

\begin{lstlisting}[style=cemdsp, caption={Diskretisierung der Ray-Hits zur Impulsantwort.}]
indices = np.round(delays * sample_rate).astype(np.int64)
ir_length = int(indices.max()) + 1
ir = np.zeros(ir_length, dtype=np.float64)
np.add.at(ir, indices, pressures)   # akkumuliert gleichzeitige Treffer
return ir
\end{lstlisting}

\subsubsection{Faltung mit \texttt{oaconvolve}}
\label{subsubsec:oaconvolve}

Die eigentliche Faltung nutzt \texttt{scipy.signal.oaconvolve}, die
Overlap-Add-basierte FFT-Faltung aus Abschnitt~2.3. Der Modus \texttt{"full"}
liefert das vollständige Ergebnis der Länge $N + M - 1$ und bewahrt damit den
kompletten Nachhall. Mehrkanaliges (etwa Stereo-)Audio wird kanalweise gefaltet
und danach wieder zusammengesetzt:

\begin{lstlisting}[style=cemdsp, caption={Mono- und Mehrkanal-Faltung.}]
if dry_audio.ndim == 1:
    wet = signal.oaconvolve(dry_audio, ir, mode="full")
elif dry_audio.ndim == 2:
    channels = [
        signal.oaconvolve(dry_audio[:, c], ir, mode="full")
        for c in range(dry_audio.shape[1])
    ]
    wet = np.stack(channels, axis=1)
return wet.astype(np.float32)
\end{lstlisting}

Vor der Faltung prüft \texttt{convolve} das Schema von \texttt{ir\_data} und
vergleicht die Abtastraten von Audio und Impulsantwort. Das Ergebnis gibt die
Funktion roh und unnormalisiert zurück. Die Lautstärke regelt erst der Export
(Abschnitt~\ref{subsec:rendering}); so bleibt die Faltung frei von Pegelfragen.

\subsubsection{Frequenzabhängige Materialdämpfung}
\label{subsubsec:material}

Reale Materialien dämpfen Schall frequenzabhängig. Ein Teppich schluckt hohe
Frequenzen viel stärker als tiefe, Beton reflektiert dagegen fast breitbandig.
Die Rust-Engine reduziert die Absorption aber schon auf einen einzigen skalaren
Restdruck $p_j$ pro Treffer, sodass die spektrale Information an dieser
Systemgrenze verloren geht. Damit der Nachhall trotzdem frequenzabhängig wird,
ergänzt \texttt{convolution.py} ein Absorptionsmodell im Frequenzbereich.

Grundlage sind die Oktavband-Absorptionskoeffizienten $\alpha_b$ aus der Datei
\texttt{materials.json}, die auch die Rust-Engine speist. Physik und
Signalverarbeitung greifen damit auf dieselbe Datenbasis zu (single source of
truth). Jedes Material ist über sechs Oktavbänder mit den Mittenfrequenzen
$125,\,250,\,500,\,1000,\,2000$ und $4000$~Hz beschrieben; ein hoher Koeffizient
bedeutet starke Dämpfung im jeweiligen Band.

Das Modell beruht auf einer einfachen Überlegung: Ein Tap, der zum Zeitpunkt $t$
eintrifft, hat näherungsweise $n(t) = t / \tau_{\mathrm{mfp}}$ Reflexionen hinter
sich, wobei $\tau_{\mathrm{mfp}}$ die mittlere Zeit zwischen zwei Reflexionen
angibt (die mittlere freie Weglänge in Zeiteinheiten). Nach $n$ Reflexionen
verbleibt in jedem Oktavband $b$ der Energieanteil
\[
    g_b(t) \;=\; \bigl(1 - \alpha_b\bigr)^{\,n(t)} \;=\; \bigl(1 - \alpha_b\bigr)^{\,t / \tau_{\mathrm{mfp}}}.
\]
Bänder mit hohem $\alpha_b$, etwa die Höhen bei Teppich, klingen dadurch deutlich
schneller ab als solche mit niedrigem $\alpha_b$. Der Nachhall wird über die Zeit
zunehmend dunkler.

Im Code setzt eine Oktavband-Filterbank dieses Modell um. Die Impulsantwort wird
über Butterworth-Bandpässe (nullphasig via \texttt{sosfiltfilt}) in ihre sechs
Bänder zerlegt, jedes Band mit der zeitabhängigen Hüllkurve $g_b(t)$ gewichtet und
die Bänder anschließend wieder aufsummiert.

\begin{lstlisting}[style=cemdsp, caption={Frequenzabhaengige Material-Daempfung der Impulsantwort.}]
centers, materials = _load_materials(str(_MATERIALS_JSON))
nyquist = sample_rate / 2.0
reflections = (np.arange(ir.size) / sample_rate) / mean_free_path_s

shaped = np.zeros_like(ir)
for (low, high), alpha in zip(_octave_band_edges(centers), materials[material]):
    band = _bandpass(ir, low, high, nyquist)      # ein Oktavband isolieren
    gain = np.power(max(1e-6, 1.0 - alpha), reflections)
    shaped += band * gain                          # zeitabhaengig daempfen
return shaped
\end{lstlisting}

\paragraph{Verifikation.} Ein reproduzierbarer Selbsttest prüft das Modell. Er
dämpft dieselbe dichte Impulsantwort einmal mit \texttt{carpet} und einmal mit
\texttt{concrete} und misst jeweils den Energieanteil oberhalb von $2$~kHz. Bei
Teppich bleiben davon nur $0{,}39$ übrig, bei Beton $0{,}71$; der Teppich schluckt
die Höhen also deutlich stärker. Kippt diese Reihenfolge, schlägt der Test fehl.

\paragraph{Grenzen des Modells.} Das Verfahren modelliert ein globales
Raummaterial, nicht die tatsächliche Materialfolge eines einzelnen Strahlpfads.
Eine Trennung nach Oberfläche, also die Aussage „dieser Tap stammt von Teppich,
jener von Beton", ist im Python-Dienst nicht möglich, solange die Rust-Engine nur
einen skalaren Druck pro Treffer liefert. Dafür müsste die Engine schon bandweise
Drücke ($p_{j,b}$ statt $p_j$) exportieren. Kapitel~7 führt das als mögliche
Erweiterung.

\subsection{Audio-Rendering und Export}
\label{subsec:rendering}

Das Ein- und Auslesen der Audiodateien sowie die Normalisierung gegen Clipping
übernimmt \texttt{wav\_handler.py}; die Details stehen in der Teildokumentation zum
Audio-I/O. Für den DSP-Teil zählt nur der Vertrag zwischen den Stufen.
\texttt{convolve} liefert ein rohes, womöglich übersteuertes Signal;
\texttt{write\_wav} misst dessen maximale Amplitude und normalisiert erst, wenn
sie den Bereich $[-1{,}0,\,1{,}0]$ verlässt. So bleibt Clipping aus, und die
Faltung muss sich nie um Pegel kümmern.

%   3. Der Master braucht nur diese Pakete (Dokumentation.tex hat beide schon):
%        \usepackage{amsmath}
%        \usepackage{listings}
%      Kein xcolor/babel noetig: Umlaute laufen unter UTF-8, deutsche
%      Anfuehrungszeichen sind literal „..." gesetzt.
%
% STANDALONE-VORSCHAU (zum alleine Kompilieren einfach drumherum setzen):
%   \documentclass{article}\usepackage{amsmath}\usepackage{listings}
%   \begin{document}
%   % =====================================================================
% Kapitel 5 - Digitale Signalverarbeitung und Orchestrierung (Python)
% Autor: Cem Acar
%
% EINBINDEN ins Hauptdokument:
%   1. Dieses File nach latexDocs/ legen.
%   2. Im Master (Dokumentation.tex) an der richtigen Stelle einfuegen:
%        \input{Kapitel5_DSP_Cem}
%   3. Der Master braucht nur diese Pakete (Dokumentation.tex hat beide schon):
%        \usepackage{amsmath}
%        \usepackage{listings}
%      Kein xcolor/babel noetig: Umlaute laufen unter UTF-8, deutsche
%      Anfuehrungszeichen sind literal „..." gesetzt.
%
% STANDALONE-VORSCHAU (zum alleine Kompilieren einfach drumherum setzen):
%   \documentclass{article}\usepackage{amsmath}\usepackage{listings}
%   \begin{document}
%   \input{Kapitel5_DSP_Cem}
%   \end{document}
% =====================================================================

% Listing-Stil fuer dieses Kapitel (nur listings-Kern, keine Farben noetig)
\lstdefinestyle{cemdsp}{
    language=Python,
    basicstyle=\ttfamily\small,
    showstringspaces=false,
    breaklines=true,
    numbers=left,
    numberstyle=\tiny,
    frame=single,
    tabsize=4
}

%======================================================================
\section{Digitale Signalverarbeitung und Orchestrierung (Python)}
\label{sec:python-dsp}

Die Rust-Engine berechnet die Physik der Schallausbreitung, der Python-Dienst
übernimmt die Signalverarbeitung und die Orchestrierung der Pipeline. Warum die
Physik in Rust und die Signalverarbeitung in Python mit NumPy und SciPy liegt,
ist in Kapitel~3 begründet. Die signaltheoretischen Grundlagen, also
Impulsantwort, Faltungsintegral und schnelle Faltung über die FFT, behandelt
Abschnitt~2.3. Dieses Kapitel zeigt, wie diese Grundlagen im Code umgesetzt sind,
und folgt dabei dem Datenfluss durch die Pipeline.

\subsection{Pipeline-Management durch \texttt{main.py}}
\label{subsec:pipeline}

\texttt{main.py} führt die vier Verarbeitungsschritte nacheinander aus. Tritt in
einem Schritt ein Fehler auf, bricht die Pipeline sofort mit einer Meldung ab,
statt mit ungültigen Zwischendaten weiterzurechnen (fail-fast).

\begin{enumerate}
    \item \textbf{Physik-Engine starten:} \texttt{run\_rust\_and\_load\_json} ruft
    die vorkompilierte Rust-Binary als Subprozess auf und lädt deren Ausgabe
    \texttt{ir\_output.json} als Python-Dictionary (\texttt{ir\_data}).
    \item \textbf{Trockenaudio laden:} \texttt{read\_wav} liest die
    Eingangsdatei und liefert Audiodaten samt Abtastrate.
    \item \textbf{Faltung:} \texttt{convolve} wendet die aus \texttt{ir\_data}
    erzeugte Impulsantwort auf das Trockensignal an.
    \item \textbf{Export:} \texttt{write\_wav} normalisiert das Ergebnis bei
    Bedarf und schreibt die fertige Ausgabedatei.
\end{enumerate}

Der Binary-Name wird plattformabhängig gewählt: unter Windows mit der Endung
\texttt{.exe}, unter macOS und Linux ohne. Stimmt die Abtastrate des Trockenaudios
nicht mit der Rate der Engine überein, gibt \texttt{main.py} eine Warnung aus,
weil der Hall sonst verstimmt klingen kann.

\subsection{Die Prozessbrücke zwischen Rust und Python}
\label{subsec:rust-runner}

Rust und Python sind nicht über eine Sprachbindung wie PyO3 gekoppelt, sondern
lose über einen Subprozess-Aufruf und eine JSON-Datei verbunden. Die Funktion
\texttt{run\_rust\_and\_load\_json} kapselt diese Brücke. Sie prüft zunächst, ob
die Binary existiert, und bestimmt das Arbeitsverzeichnis anhand der
\texttt{input\_config.json}. Eine alte Ausgabedatei wird vorher gelöscht, damit
kein veraltetes Ergebnis eingelesen wird. Nach dem Lauf wertet die Funktion den
Rückgabecode aus; schlägt Rust fehl, landen \texttt{stdout} und \texttt{stderr} in
der Fehlermeldung und erleichtern die Diagnose. Der Vorteil der losen Kopplung
liegt im verteilten Betrieb: Beide Dienste lassen sich getrennt kompilieren,
testen und auf mehrere Cluster-Knoten verteilen.

\subsection{Signalverarbeitung mit SciPy}
\label{subsec:scipy}

Das Modul \texttt{convolution.py} bildet den Kern des Python-Dienstes. Es wandelt
die rohen Verzögerungs- und Druckdaten in eine abtastgenaue Impulsantwort um und
faltet damit das Trockensignal.

\subsubsection{Diskretisierung der Impulsantwort}
\label{subsubsec:build-ir}

Die Rust-Engine liefert eine endliche Menge diskreter Strahl-Treffer, jeweils
bestehend aus einer Ankunftsverzögerung $\tau_j$ (in Sekunden) und einem skalaren
Restdruck $p_j$. Die Funktion \texttt{build\_impulse\_response} quantisiert jeden
Treffer auf den nächstgelegenen Abtastindex $i_j = \operatorname{round}(\tau_j
\cdot f_s)$ und akkumuliert gleichzeitige Treffer zur diskreten Impulsantwort:
\[
    h[i] \;=\; \sum_{j} p_j \cdot \bigl[\,i_j = i\,\bigr].
\]
Die Akkumulation erfolgt mit \texttt{numpy.add.at}. Anders als die einfache
indizierte Zuweisung, die bei doppelten Indizes nur den letzten Wert behält,
summiert \texttt{add.at} gleichzeitige Treffer korrekt auf. Vor der Berechnung
prüft die Funktion die Eingaben an der Systemgrenze: übereinstimmende
Array-Formen, Eindimensionalität, nicht-negative Verzögerungen und eine positive
Abtastrate.

\begin{lstlisting}[style=cemdsp, caption={Diskretisierung der Ray-Hits zur Impulsantwort.}]
indices = np.round(delays * sample_rate).astype(np.int64)
ir_length = int(indices.max()) + 1
ir = np.zeros(ir_length, dtype=np.float64)
np.add.at(ir, indices, pressures)   # akkumuliert gleichzeitige Treffer
return ir
\end{lstlisting}

\subsubsection{Faltung mit \texttt{oaconvolve}}
\label{subsubsec:oaconvolve}

Die eigentliche Faltung nutzt \texttt{scipy.signal.oaconvolve}, die
Overlap-Add-basierte FFT-Faltung aus Abschnitt~2.3. Der Modus \texttt{"full"}
liefert das vollständige Ergebnis der Länge $N + M - 1$ und bewahrt damit den
kompletten Nachhall. Mehrkanaliges (etwa Stereo-)Audio wird kanalweise gefaltet
und danach wieder zusammengesetzt:

\begin{lstlisting}[style=cemdsp, caption={Mono- und Mehrkanal-Faltung.}]
if dry_audio.ndim == 1:
    wet = signal.oaconvolve(dry_audio, ir, mode="full")
elif dry_audio.ndim == 2:
    channels = [
        signal.oaconvolve(dry_audio[:, c], ir, mode="full")
        for c in range(dry_audio.shape[1])
    ]
    wet = np.stack(channels, axis=1)
return wet.astype(np.float32)
\end{lstlisting}

Vor der Faltung prüft \texttt{convolve} das Schema von \texttt{ir\_data} und
vergleicht die Abtastraten von Audio und Impulsantwort. Das Ergebnis gibt die
Funktion roh und unnormalisiert zurück. Die Lautstärke regelt erst der Export
(Abschnitt~\ref{subsec:rendering}); so bleibt die Faltung frei von Pegelfragen.

\subsubsection{Frequenzabhängige Materialdämpfung}
\label{subsubsec:material}

Reale Materialien dämpfen Schall frequenzabhängig. Ein Teppich schluckt hohe
Frequenzen viel stärker als tiefe, Beton reflektiert dagegen fast breitbandig.
Die Rust-Engine reduziert die Absorption aber schon auf einen einzigen skalaren
Restdruck $p_j$ pro Treffer, sodass die spektrale Information an dieser
Systemgrenze verloren geht. Damit der Nachhall trotzdem frequenzabhängig wird,
ergänzt \texttt{convolution.py} ein Absorptionsmodell im Frequenzbereich.

Grundlage sind die Oktavband-Absorptionskoeffizienten $\alpha_b$ aus der Datei
\texttt{materials.json}, die auch die Rust-Engine speist. Physik und
Signalverarbeitung greifen damit auf dieselbe Datenbasis zu (single source of
truth). Jedes Material ist über sechs Oktavbänder mit den Mittenfrequenzen
$125,\,250,\,500,\,1000,\,2000$ und $4000$~Hz beschrieben; ein hoher Koeffizient
bedeutet starke Dämpfung im jeweiligen Band.

Das Modell beruht auf einer einfachen Überlegung: Ein Tap, der zum Zeitpunkt $t$
eintrifft, hat näherungsweise $n(t) = t / \tau_{\mathrm{mfp}}$ Reflexionen hinter
sich, wobei $\tau_{\mathrm{mfp}}$ die mittlere Zeit zwischen zwei Reflexionen
angibt (die mittlere freie Weglänge in Zeiteinheiten). Nach $n$ Reflexionen
verbleibt in jedem Oktavband $b$ der Energieanteil
\[
    g_b(t) \;=\; \bigl(1 - \alpha_b\bigr)^{\,n(t)} \;=\; \bigl(1 - \alpha_b\bigr)^{\,t / \tau_{\mathrm{mfp}}}.
\]
Bänder mit hohem $\alpha_b$, etwa die Höhen bei Teppich, klingen dadurch deutlich
schneller ab als solche mit niedrigem $\alpha_b$. Der Nachhall wird über die Zeit
zunehmend dunkler.

Im Code setzt eine Oktavband-Filterbank dieses Modell um. Die Impulsantwort wird
über Butterworth-Bandpässe (nullphasig via \texttt{sosfiltfilt}) in ihre sechs
Bänder zerlegt, jedes Band mit der zeitabhängigen Hüllkurve $g_b(t)$ gewichtet und
die Bänder anschließend wieder aufsummiert.

\begin{lstlisting}[style=cemdsp, caption={Frequenzabhaengige Material-Daempfung der Impulsantwort.}]
centers, materials = _load_materials(str(_MATERIALS_JSON))
nyquist = sample_rate / 2.0
reflections = (np.arange(ir.size) / sample_rate) / mean_free_path_s

shaped = np.zeros_like(ir)
for (low, high), alpha in zip(_octave_band_edges(centers), materials[material]):
    band = _bandpass(ir, low, high, nyquist)      # ein Oktavband isolieren
    gain = np.power(max(1e-6, 1.0 - alpha), reflections)
    shaped += band * gain                          # zeitabhaengig daempfen
return shaped
\end{lstlisting}

\paragraph{Verifikation.} Ein reproduzierbarer Selbsttest prüft das Modell. Er
dämpft dieselbe dichte Impulsantwort einmal mit \texttt{carpet} und einmal mit
\texttt{concrete} und misst jeweils den Energieanteil oberhalb von $2$~kHz. Bei
Teppich bleiben davon nur $0{,}39$ übrig, bei Beton $0{,}71$; der Teppich schluckt
die Höhen also deutlich stärker. Kippt diese Reihenfolge, schlägt der Test fehl.

\paragraph{Grenzen des Modells.} Das Verfahren modelliert ein globales
Raummaterial, nicht die tatsächliche Materialfolge eines einzelnen Strahlpfads.
Eine Trennung nach Oberfläche, also die Aussage „dieser Tap stammt von Teppich,
jener von Beton", ist im Python-Dienst nicht möglich, solange die Rust-Engine nur
einen skalaren Druck pro Treffer liefert. Dafür müsste die Engine schon bandweise
Drücke ($p_{j,b}$ statt $p_j$) exportieren. Kapitel~7 führt das als mögliche
Erweiterung.

\subsection{Audio-Rendering und Export}
\label{subsec:rendering}

Das Ein- und Auslesen der Audiodateien sowie die Normalisierung gegen Clipping
übernimmt \texttt{wav\_handler.py}; die Details stehen in der Teildokumentation zum
Audio-I/O. Für den DSP-Teil zählt nur der Vertrag zwischen den Stufen.
\texttt{convolve} liefert ein rohes, womöglich übersteuertes Signal;
\texttt{write\_wav} misst dessen maximale Amplitude und normalisiert erst, wenn
sie den Bereich $[-1{,}0,\,1{,}0]$ verlässt. So bleibt Clipping aus, und die
Faltung muss sich nie um Pegel kümmern.

%   \end{document}
% =====================================================================

% Listing-Stil fuer dieses Kapitel (nur listings-Kern, keine Farben noetig)
\lstdefinestyle{cemdsp}{
    language=Python,
    basicstyle=\ttfamily\small,
    showstringspaces=false,
    breaklines=true,
    numbers=left,
    numberstyle=\tiny,
    frame=single,
    tabsize=4
}

%======================================================================
\section{Digitale Signalverarbeitung und Orchestrierung (Python)}
\label{sec:python-dsp}

Die Rust-Engine berechnet die Physik der Schallausbreitung, der Python-Dienst
übernimmt die Signalverarbeitung und die Orchestrierung der Pipeline. Warum die
Physik in Rust und die Signalverarbeitung in Python mit NumPy und SciPy liegt,
ist in Kapitel~3 begründet. Die signaltheoretischen Grundlagen, also
Impulsantwort, Faltungsintegral und schnelle Faltung über die FFT, behandelt
Abschnitt~2.3. Dieses Kapitel zeigt, wie diese Grundlagen im Code umgesetzt sind,
und folgt dabei dem Datenfluss durch die Pipeline.

\subsection{Pipeline-Management durch \texttt{main.py}}
\label{subsec:pipeline}

\texttt{main.py} führt die vier Verarbeitungsschritte nacheinander aus. Tritt in
einem Schritt ein Fehler auf, bricht die Pipeline sofort mit einer Meldung ab,
statt mit ungültigen Zwischendaten weiterzurechnen (fail-fast).

\begin{enumerate}
    \item \textbf{Physik-Engine starten:} \texttt{run\_rust\_and\_load\_json} ruft
    die vorkompilierte Rust-Binary als Subprozess auf und lädt deren Ausgabe
    \texttt{ir\_output.json} als Python-Dictionary (\texttt{ir\_data}).
    \item \textbf{Trockenaudio laden:} \texttt{read\_wav} liest die
    Eingangsdatei und liefert Audiodaten samt Abtastrate.
    \item \textbf{Faltung:} \texttt{convolve} wendet die aus \texttt{ir\_data}
    erzeugte Impulsantwort auf das Trockensignal an.
    \item \textbf{Export:} \texttt{write\_wav} normalisiert das Ergebnis bei
    Bedarf und schreibt die fertige Ausgabedatei.
\end{enumerate}

Der Binary-Name wird plattformabhängig gewählt: unter Windows mit der Endung
\texttt{.exe}, unter macOS und Linux ohne. Stimmt die Abtastrate des Trockenaudios
nicht mit der Rate der Engine überein, gibt \texttt{main.py} eine Warnung aus,
weil der Hall sonst verstimmt klingen kann.

\subsection{Die Prozessbrücke zwischen Rust und Python}
\label{subsec:rust-runner}

Rust und Python sind nicht über eine Sprachbindung wie PyO3 gekoppelt, sondern
lose über einen Subprozess-Aufruf und eine JSON-Datei verbunden. Die Funktion
\texttt{run\_rust\_and\_load\_json} kapselt diese Brücke. Sie prüft zunächst, ob
die Binary existiert, und bestimmt das Arbeitsverzeichnis anhand der
\texttt{input\_config.json}. Eine alte Ausgabedatei wird vorher gelöscht, damit
kein veraltetes Ergebnis eingelesen wird. Nach dem Lauf wertet die Funktion den
Rückgabecode aus; schlägt Rust fehl, landen \texttt{stdout} und \texttt{stderr} in
der Fehlermeldung und erleichtern die Diagnose. Der Vorteil der losen Kopplung
liegt im verteilten Betrieb: Beide Dienste lassen sich getrennt kompilieren,
testen und auf mehrere Cluster-Knoten verteilen.

\subsection{Signalverarbeitung mit SciPy}
\label{subsec:scipy}

Das Modul \texttt{convolution.py} bildet den Kern des Python-Dienstes. Es wandelt
die rohen Verzögerungs- und Druckdaten in eine abtastgenaue Impulsantwort um und
faltet damit das Trockensignal.

\subsubsection{Diskretisierung der Impulsantwort}
\label{subsubsec:build-ir}

Die Rust-Engine liefert eine endliche Menge diskreter Strahl-Treffer, jeweils
bestehend aus einer Ankunftsverzögerung $\tau_j$ (in Sekunden) und einem skalaren
Restdruck $p_j$. Die Funktion \texttt{build\_impulse\_response} quantisiert jeden
Treffer auf den nächstgelegenen Abtastindex $i_j = \operatorname{round}(\tau_j
\cdot f_s)$ und akkumuliert gleichzeitige Treffer zur diskreten Impulsantwort:
\[
    h[i] \;=\; \sum_{j} p_j \cdot \bigl[\,i_j = i\,\bigr].
\]
Die Akkumulation erfolgt mit \texttt{numpy.add.at}. Anders als die einfache
indizierte Zuweisung, die bei doppelten Indizes nur den letzten Wert behält,
summiert \texttt{add.at} gleichzeitige Treffer korrekt auf. Vor der Berechnung
prüft die Funktion die Eingaben an der Systemgrenze: übereinstimmende
Array-Formen, Eindimensionalität, nicht-negative Verzögerungen und eine positive
Abtastrate.

\begin{lstlisting}[style=cemdsp, caption={Diskretisierung der Ray-Hits zur Impulsantwort.}]
indices = np.round(delays * sample_rate).astype(np.int64)
ir_length = int(indices.max()) + 1
ir = np.zeros(ir_length, dtype=np.float64)
np.add.at(ir, indices, pressures)   # akkumuliert gleichzeitige Treffer
return ir
\end{lstlisting}

\subsubsection{Faltung mit \texttt{oaconvolve}}
\label{subsubsec:oaconvolve}

Die eigentliche Faltung nutzt \texttt{scipy.signal.oaconvolve}, die
Overlap-Add-basierte FFT-Faltung aus Abschnitt~2.3. Der Modus \texttt{"full"}
liefert das vollständige Ergebnis der Länge $N + M - 1$ und bewahrt damit den
kompletten Nachhall. Mehrkanaliges (etwa Stereo-)Audio wird kanalweise gefaltet
und danach wieder zusammengesetzt:

\begin{lstlisting}[style=cemdsp, caption={Mono- und Mehrkanal-Faltung.}]
if dry_audio.ndim == 1:
    wet = signal.oaconvolve(dry_audio, ir, mode="full")
elif dry_audio.ndim == 2:
    channels = [
        signal.oaconvolve(dry_audio[:, c], ir, mode="full")
        for c in range(dry_audio.shape[1])
    ]
    wet = np.stack(channels, axis=1)
return wet.astype(np.float32)
\end{lstlisting}

Vor der Faltung prüft \texttt{convolve} das Schema von \texttt{ir\_data} und
vergleicht die Abtastraten von Audio und Impulsantwort. Das Ergebnis gibt die
Funktion roh und unnormalisiert zurück. Die Lautstärke regelt erst der Export
(Abschnitt~\ref{subsec:rendering}); so bleibt die Faltung frei von Pegelfragen.

\subsubsection{Frequenzabhängige Materialdämpfung}
\label{subsubsec:material}

Reale Materialien dämpfen Schall frequenzabhängig. Ein Teppich schluckt hohe
Frequenzen viel stärker als tiefe, Beton reflektiert dagegen fast breitbandig.
Die Rust-Engine reduziert die Absorption aber schon auf einen einzigen skalaren
Restdruck $p_j$ pro Treffer, sodass die spektrale Information an dieser
Systemgrenze verloren geht. Damit der Nachhall trotzdem frequenzabhängig wird,
ergänzt \texttt{convolution.py} ein Absorptionsmodell im Frequenzbereich.

Grundlage sind die Oktavband-Absorptionskoeffizienten $\alpha_b$ aus der Datei
\texttt{materials.json}, die auch die Rust-Engine speist. Physik und
Signalverarbeitung greifen damit auf dieselbe Datenbasis zu (single source of
truth). Jedes Material ist über sechs Oktavbänder mit den Mittenfrequenzen
$125,\,250,\,500,\,1000,\,2000$ und $4000$~Hz beschrieben; ein hoher Koeffizient
bedeutet starke Dämpfung im jeweiligen Band.

Das Modell beruht auf einer einfachen Überlegung: Ein Tap, der zum Zeitpunkt $t$
eintrifft, hat näherungsweise $n(t) = t / \tau_{\mathrm{mfp}}$ Reflexionen hinter
sich, wobei $\tau_{\mathrm{mfp}}$ die mittlere Zeit zwischen zwei Reflexionen
angibt (die mittlere freie Weglänge in Zeiteinheiten). Nach $n$ Reflexionen
verbleibt in jedem Oktavband $b$ der Energieanteil
\[
    g_b(t) \;=\; \bigl(1 - \alpha_b\bigr)^{\,n(t)} \;=\; \bigl(1 - \alpha_b\bigr)^{\,t / \tau_{\mathrm{mfp}}}.
\]
Bänder mit hohem $\alpha_b$, etwa die Höhen bei Teppich, klingen dadurch deutlich
schneller ab als solche mit niedrigem $\alpha_b$. Der Nachhall wird über die Zeit
zunehmend dunkler.

Im Code setzt eine Oktavband-Filterbank dieses Modell um. Die Impulsantwort wird
über Butterworth-Bandpässe (nullphasig via \texttt{sosfiltfilt}) in ihre sechs
Bänder zerlegt, jedes Band mit der zeitabhängigen Hüllkurve $g_b(t)$ gewichtet und
die Bänder anschließend wieder aufsummiert.

\begin{lstlisting}[style=cemdsp, caption={Frequenzabhaengige Material-Daempfung der Impulsantwort.}]
centers, materials = _load_materials(str(_MATERIALS_JSON))
nyquist = sample_rate / 2.0
reflections = (np.arange(ir.size) / sample_rate) / mean_free_path_s

shaped = np.zeros_like(ir)
for (low, high), alpha in zip(_octave_band_edges(centers), materials[material]):
    band = _bandpass(ir, low, high, nyquist)      # ein Oktavband isolieren
    gain = np.power(max(1e-6, 1.0 - alpha), reflections)
    shaped += band * gain                          # zeitabhaengig daempfen
return shaped
\end{lstlisting}

\paragraph{Verifikation.} Ein reproduzierbarer Selbsttest prüft das Modell. Er
dämpft dieselbe dichte Impulsantwort einmal mit \texttt{carpet} und einmal mit
\texttt{concrete} und misst jeweils den Energieanteil oberhalb von $2$~kHz. Bei
Teppich bleiben davon nur $0{,}39$ übrig, bei Beton $0{,}71$; der Teppich schluckt
die Höhen also deutlich stärker. Kippt diese Reihenfolge, schlägt der Test fehl.

\paragraph{Grenzen des Modells.} Das Verfahren modelliert ein globales
Raummaterial, nicht die tatsächliche Materialfolge eines einzelnen Strahlpfads.
Eine Trennung nach Oberfläche, also die Aussage „dieser Tap stammt von Teppich,
jener von Beton", ist im Python-Dienst nicht möglich, solange die Rust-Engine nur
einen skalaren Druck pro Treffer liefert. Dafür müsste die Engine schon bandweise
Drücke ($p_{j,b}$ statt $p_j$) exportieren. Kapitel~7 führt das als mögliche
Erweiterung.

\subsection{Audio-Rendering und Export}
\label{subsec:rendering}

Das Ein- und Auslesen der Audiodateien sowie die Normalisierung gegen Clipping
übernimmt \texttt{wav\_handler.py}; die Details stehen in der Teildokumentation zum
Audio-I/O. Für den DSP-Teil zählt nur der Vertrag zwischen den Stufen.
\texttt{convolve} liefert ein rohes, womöglich übersteuertes Signal;
\texttt{write\_wav} misst dessen maximale Amplitude und normalisiert erst, wenn
sie den Bereich $[-1{,}0,\,1{,}0]$ verlässt. So bleibt Clipping aus, und die
Faltung muss sich nie um Pegel kümmern.

%   3. Der Master braucht nur diese Pakete (Dokumentation.tex hat beide schon):
%        \usepackage{amsmath}
%        \usepackage{listings}
%      Kein xcolor/babel noetig: Umlaute laufen unter UTF-8, deutsche
%      Anfuehrungszeichen sind literal „..." gesetzt.
%
% STANDALONE-VORSCHAU (zum alleine Kompilieren einfach drumherum setzen):
%   \documentclass{article}\usepackage{amsmath}\usepackage{listings}
%   \begin{document}
%   % =====================================================================
% Kapitel 5 - Digitale Signalverarbeitung und Orchestrierung (Python)
% Autor: Cem Acar
%
% EINBINDEN ins Hauptdokument:
%   1. Dieses File nach latexDocs/ legen.
%   2. Im Master (Dokumentation.tex) an der richtigen Stelle einfuegen:
%        % =====================================================================
% Kapitel 5 - Digitale Signalverarbeitung und Orchestrierung (Python)
% Autor: Cem Acar
%
% EINBINDEN ins Hauptdokument:
%   1. Dieses File nach latexDocs/ legen.
%   2. Im Master (Dokumentation.tex) an der richtigen Stelle einfuegen:
%        \input{Kapitel5_DSP_Cem}
%   3. Der Master braucht nur diese Pakete (Dokumentation.tex hat beide schon):
%        \usepackage{amsmath}
%        \usepackage{listings}
%      Kein xcolor/babel noetig: Umlaute laufen unter UTF-8, deutsche
%      Anfuehrungszeichen sind literal „..." gesetzt.
%
% STANDALONE-VORSCHAU (zum alleine Kompilieren einfach drumherum setzen):
%   \documentclass{article}\usepackage{amsmath}\usepackage{listings}
%   \begin{document}
%   \input{Kapitel5_DSP_Cem}
%   \end{document}
% =====================================================================

% Listing-Stil fuer dieses Kapitel (nur listings-Kern, keine Farben noetig)
\lstdefinestyle{cemdsp}{
    language=Python,
    basicstyle=\ttfamily\small,
    showstringspaces=false,
    breaklines=true,
    numbers=left,
    numberstyle=\tiny,
    frame=single,
    tabsize=4
}

%======================================================================
\section{Digitale Signalverarbeitung und Orchestrierung (Python)}
\label{sec:python-dsp}

Die Rust-Engine berechnet die Physik der Schallausbreitung, der Python-Dienst
übernimmt die Signalverarbeitung und die Orchestrierung der Pipeline. Warum die
Physik in Rust und die Signalverarbeitung in Python mit NumPy und SciPy liegt,
ist in Kapitel~3 begründet. Die signaltheoretischen Grundlagen, also
Impulsantwort, Faltungsintegral und schnelle Faltung über die FFT, behandelt
Abschnitt~2.3. Dieses Kapitel zeigt, wie diese Grundlagen im Code umgesetzt sind,
und folgt dabei dem Datenfluss durch die Pipeline.

\subsection{Pipeline-Management durch \texttt{main.py}}
\label{subsec:pipeline}

\texttt{main.py} führt die vier Verarbeitungsschritte nacheinander aus. Tritt in
einem Schritt ein Fehler auf, bricht die Pipeline sofort mit einer Meldung ab,
statt mit ungültigen Zwischendaten weiterzurechnen (fail-fast).

\begin{enumerate}
    \item \textbf{Physik-Engine starten:} \texttt{run\_rust\_and\_load\_json} ruft
    die vorkompilierte Rust-Binary als Subprozess auf und lädt deren Ausgabe
    \texttt{ir\_output.json} als Python-Dictionary (\texttt{ir\_data}).
    \item \textbf{Trockenaudio laden:} \texttt{read\_wav} liest die
    Eingangsdatei und liefert Audiodaten samt Abtastrate.
    \item \textbf{Faltung:} \texttt{convolve} wendet die aus \texttt{ir\_data}
    erzeugte Impulsantwort auf das Trockensignal an.
    \item \textbf{Export:} \texttt{write\_wav} normalisiert das Ergebnis bei
    Bedarf und schreibt die fertige Ausgabedatei.
\end{enumerate}

Der Binary-Name wird plattformabhängig gewählt: unter Windows mit der Endung
\texttt{.exe}, unter macOS und Linux ohne. Stimmt die Abtastrate des Trockenaudios
nicht mit der Rate der Engine überein, gibt \texttt{main.py} eine Warnung aus,
weil der Hall sonst verstimmt klingen kann.

\subsection{Die Prozessbrücke zwischen Rust und Python}
\label{subsec:rust-runner}

Rust und Python sind nicht über eine Sprachbindung wie PyO3 gekoppelt, sondern
lose über einen Subprozess-Aufruf und eine JSON-Datei verbunden. Die Funktion
\texttt{run\_rust\_and\_load\_json} kapselt diese Brücke. Sie prüft zunächst, ob
die Binary existiert, und bestimmt das Arbeitsverzeichnis anhand der
\texttt{input\_config.json}. Eine alte Ausgabedatei wird vorher gelöscht, damit
kein veraltetes Ergebnis eingelesen wird. Nach dem Lauf wertet die Funktion den
Rückgabecode aus; schlägt Rust fehl, landen \texttt{stdout} und \texttt{stderr} in
der Fehlermeldung und erleichtern die Diagnose. Der Vorteil der losen Kopplung
liegt im verteilten Betrieb: Beide Dienste lassen sich getrennt kompilieren,
testen und auf mehrere Cluster-Knoten verteilen.

\subsection{Signalverarbeitung mit SciPy}
\label{subsec:scipy}

Das Modul \texttt{convolution.py} bildet den Kern des Python-Dienstes. Es wandelt
die rohen Verzögerungs- und Druckdaten in eine abtastgenaue Impulsantwort um und
faltet damit das Trockensignal.

\subsubsection{Diskretisierung der Impulsantwort}
\label{subsubsec:build-ir}

Die Rust-Engine liefert eine endliche Menge diskreter Strahl-Treffer, jeweils
bestehend aus einer Ankunftsverzögerung $\tau_j$ (in Sekunden) und einem skalaren
Restdruck $p_j$. Die Funktion \texttt{build\_impulse\_response} quantisiert jeden
Treffer auf den nächstgelegenen Abtastindex $i_j = \operatorname{round}(\tau_j
\cdot f_s)$ und akkumuliert gleichzeitige Treffer zur diskreten Impulsantwort:
\[
    h[i] \;=\; \sum_{j} p_j \cdot \bigl[\,i_j = i\,\bigr].
\]
Die Akkumulation erfolgt mit \texttt{numpy.add.at}. Anders als die einfache
indizierte Zuweisung, die bei doppelten Indizes nur den letzten Wert behält,
summiert \texttt{add.at} gleichzeitige Treffer korrekt auf. Vor der Berechnung
prüft die Funktion die Eingaben an der Systemgrenze: übereinstimmende
Array-Formen, Eindimensionalität, nicht-negative Verzögerungen und eine positive
Abtastrate.

\begin{lstlisting}[style=cemdsp, caption={Diskretisierung der Ray-Hits zur Impulsantwort.}]
indices = np.round(delays * sample_rate).astype(np.int64)
ir_length = int(indices.max()) + 1
ir = np.zeros(ir_length, dtype=np.float64)
np.add.at(ir, indices, pressures)   # akkumuliert gleichzeitige Treffer
return ir
\end{lstlisting}

\subsubsection{Faltung mit \texttt{oaconvolve}}
\label{subsubsec:oaconvolve}

Die eigentliche Faltung nutzt \texttt{scipy.signal.oaconvolve}, die
Overlap-Add-basierte FFT-Faltung aus Abschnitt~2.3. Der Modus \texttt{"full"}
liefert das vollständige Ergebnis der Länge $N + M - 1$ und bewahrt damit den
kompletten Nachhall. Mehrkanaliges (etwa Stereo-)Audio wird kanalweise gefaltet
und danach wieder zusammengesetzt:

\begin{lstlisting}[style=cemdsp, caption={Mono- und Mehrkanal-Faltung.}]
if dry_audio.ndim == 1:
    wet = signal.oaconvolve(dry_audio, ir, mode="full")
elif dry_audio.ndim == 2:
    channels = [
        signal.oaconvolve(dry_audio[:, c], ir, mode="full")
        for c in range(dry_audio.shape[1])
    ]
    wet = np.stack(channels, axis=1)
return wet.astype(np.float32)
\end{lstlisting}

Vor der Faltung prüft \texttt{convolve} das Schema von \texttt{ir\_data} und
vergleicht die Abtastraten von Audio und Impulsantwort. Das Ergebnis gibt die
Funktion roh und unnormalisiert zurück. Die Lautstärke regelt erst der Export
(Abschnitt~\ref{subsec:rendering}); so bleibt die Faltung frei von Pegelfragen.

\subsubsection{Frequenzabhängige Materialdämpfung}
\label{subsubsec:material}

Reale Materialien dämpfen Schall frequenzabhängig. Ein Teppich schluckt hohe
Frequenzen viel stärker als tiefe, Beton reflektiert dagegen fast breitbandig.
Die Rust-Engine reduziert die Absorption aber schon auf einen einzigen skalaren
Restdruck $p_j$ pro Treffer, sodass die spektrale Information an dieser
Systemgrenze verloren geht. Damit der Nachhall trotzdem frequenzabhängig wird,
ergänzt \texttt{convolution.py} ein Absorptionsmodell im Frequenzbereich.

Grundlage sind die Oktavband-Absorptionskoeffizienten $\alpha_b$ aus der Datei
\texttt{materials.json}, die auch die Rust-Engine speist. Physik und
Signalverarbeitung greifen damit auf dieselbe Datenbasis zu (single source of
truth). Jedes Material ist über sechs Oktavbänder mit den Mittenfrequenzen
$125,\,250,\,500,\,1000,\,2000$ und $4000$~Hz beschrieben; ein hoher Koeffizient
bedeutet starke Dämpfung im jeweiligen Band.

Das Modell beruht auf einer einfachen Überlegung: Ein Tap, der zum Zeitpunkt $t$
eintrifft, hat näherungsweise $n(t) = t / \tau_{\mathrm{mfp}}$ Reflexionen hinter
sich, wobei $\tau_{\mathrm{mfp}}$ die mittlere Zeit zwischen zwei Reflexionen
angibt (die mittlere freie Weglänge in Zeiteinheiten). Nach $n$ Reflexionen
verbleibt in jedem Oktavband $b$ der Energieanteil
\[
    g_b(t) \;=\; \bigl(1 - \alpha_b\bigr)^{\,n(t)} \;=\; \bigl(1 - \alpha_b\bigr)^{\,t / \tau_{\mathrm{mfp}}}.
\]
Bänder mit hohem $\alpha_b$, etwa die Höhen bei Teppich, klingen dadurch deutlich
schneller ab als solche mit niedrigem $\alpha_b$. Der Nachhall wird über die Zeit
zunehmend dunkler.

Im Code setzt eine Oktavband-Filterbank dieses Modell um. Die Impulsantwort wird
über Butterworth-Bandpässe (nullphasig via \texttt{sosfiltfilt}) in ihre sechs
Bänder zerlegt, jedes Band mit der zeitabhängigen Hüllkurve $g_b(t)$ gewichtet und
die Bänder anschließend wieder aufsummiert.

\begin{lstlisting}[style=cemdsp, caption={Frequenzabhaengige Material-Daempfung der Impulsantwort.}]
centers, materials = _load_materials(str(_MATERIALS_JSON))
nyquist = sample_rate / 2.0
reflections = (np.arange(ir.size) / sample_rate) / mean_free_path_s

shaped = np.zeros_like(ir)
for (low, high), alpha in zip(_octave_band_edges(centers), materials[material]):
    band = _bandpass(ir, low, high, nyquist)      # ein Oktavband isolieren
    gain = np.power(max(1e-6, 1.0 - alpha), reflections)
    shaped += band * gain                          # zeitabhaengig daempfen
return shaped
\end{lstlisting}

\paragraph{Verifikation.} Ein reproduzierbarer Selbsttest prüft das Modell. Er
dämpft dieselbe dichte Impulsantwort einmal mit \texttt{carpet} und einmal mit
\texttt{concrete} und misst jeweils den Energieanteil oberhalb von $2$~kHz. Bei
Teppich bleiben davon nur $0{,}39$ übrig, bei Beton $0{,}71$; der Teppich schluckt
die Höhen also deutlich stärker. Kippt diese Reihenfolge, schlägt der Test fehl.

\paragraph{Grenzen des Modells.} Das Verfahren modelliert ein globales
Raummaterial, nicht die tatsächliche Materialfolge eines einzelnen Strahlpfads.
Eine Trennung nach Oberfläche, also die Aussage „dieser Tap stammt von Teppich,
jener von Beton", ist im Python-Dienst nicht möglich, solange die Rust-Engine nur
einen skalaren Druck pro Treffer liefert. Dafür müsste die Engine schon bandweise
Drücke ($p_{j,b}$ statt $p_j$) exportieren. Kapitel~7 führt das als mögliche
Erweiterung.

\subsection{Audio-Rendering und Export}
\label{subsec:rendering}

Das Ein- und Auslesen der Audiodateien sowie die Normalisierung gegen Clipping
übernimmt \texttt{wav\_handler.py}; die Details stehen in der Teildokumentation zum
Audio-I/O. Für den DSP-Teil zählt nur der Vertrag zwischen den Stufen.
\texttt{convolve} liefert ein rohes, womöglich übersteuertes Signal;
\texttt{write\_wav} misst dessen maximale Amplitude und normalisiert erst, wenn
sie den Bereich $[-1{,}0,\,1{,}0]$ verlässt. So bleibt Clipping aus, und die
Faltung muss sich nie um Pegel kümmern.

%   3. Der Master braucht nur diese Pakete (Dokumentation.tex hat beide schon):
%        \usepackage{amsmath}
%        \usepackage{listings}
%      Kein xcolor/babel noetig: Umlaute laufen unter UTF-8, deutsche
%      Anfuehrungszeichen sind literal „..." gesetzt.
%
% STANDALONE-VORSCHAU (zum alleine Kompilieren einfach drumherum setzen):
%   \documentclass{article}\usepackage{amsmath}\usepackage{listings}
%   \begin{document}
%   % =====================================================================
% Kapitel 5 - Digitale Signalverarbeitung und Orchestrierung (Python)
% Autor: Cem Acar
%
% EINBINDEN ins Hauptdokument:
%   1. Dieses File nach latexDocs/ legen.
%   2. Im Master (Dokumentation.tex) an der richtigen Stelle einfuegen:
%        \input{Kapitel5_DSP_Cem}
%   3. Der Master braucht nur diese Pakete (Dokumentation.tex hat beide schon):
%        \usepackage{amsmath}
%        \usepackage{listings}
%      Kein xcolor/babel noetig: Umlaute laufen unter UTF-8, deutsche
%      Anfuehrungszeichen sind literal „..." gesetzt.
%
% STANDALONE-VORSCHAU (zum alleine Kompilieren einfach drumherum setzen):
%   \documentclass{article}\usepackage{amsmath}\usepackage{listings}
%   \begin{document}
%   \input{Kapitel5_DSP_Cem}
%   \end{document}
% =====================================================================

% Listing-Stil fuer dieses Kapitel (nur listings-Kern, keine Farben noetig)
\lstdefinestyle{cemdsp}{
    language=Python,
    basicstyle=\ttfamily\small,
    showstringspaces=false,
    breaklines=true,
    numbers=left,
    numberstyle=\tiny,
    frame=single,
    tabsize=4
}

%======================================================================
\section{Digitale Signalverarbeitung und Orchestrierung (Python)}
\label{sec:python-dsp}

Die Rust-Engine berechnet die Physik der Schallausbreitung, der Python-Dienst
übernimmt die Signalverarbeitung und die Orchestrierung der Pipeline. Warum die
Physik in Rust und die Signalverarbeitung in Python mit NumPy und SciPy liegt,
ist in Kapitel~3 begründet. Die signaltheoretischen Grundlagen, also
Impulsantwort, Faltungsintegral und schnelle Faltung über die FFT, behandelt
Abschnitt~2.3. Dieses Kapitel zeigt, wie diese Grundlagen im Code umgesetzt sind,
und folgt dabei dem Datenfluss durch die Pipeline.

\subsection{Pipeline-Management durch \texttt{main.py}}
\label{subsec:pipeline}

\texttt{main.py} führt die vier Verarbeitungsschritte nacheinander aus. Tritt in
einem Schritt ein Fehler auf, bricht die Pipeline sofort mit einer Meldung ab,
statt mit ungültigen Zwischendaten weiterzurechnen (fail-fast).

\begin{enumerate}
    \item \textbf{Physik-Engine starten:} \texttt{run\_rust\_and\_load\_json} ruft
    die vorkompilierte Rust-Binary als Subprozess auf und lädt deren Ausgabe
    \texttt{ir\_output.json} als Python-Dictionary (\texttt{ir\_data}).
    \item \textbf{Trockenaudio laden:} \texttt{read\_wav} liest die
    Eingangsdatei und liefert Audiodaten samt Abtastrate.
    \item \textbf{Faltung:} \texttt{convolve} wendet die aus \texttt{ir\_data}
    erzeugte Impulsantwort auf das Trockensignal an.
    \item \textbf{Export:} \texttt{write\_wav} normalisiert das Ergebnis bei
    Bedarf und schreibt die fertige Ausgabedatei.
\end{enumerate}

Der Binary-Name wird plattformabhängig gewählt: unter Windows mit der Endung
\texttt{.exe}, unter macOS und Linux ohne. Stimmt die Abtastrate des Trockenaudios
nicht mit der Rate der Engine überein, gibt \texttt{main.py} eine Warnung aus,
weil der Hall sonst verstimmt klingen kann.

\subsection{Die Prozessbrücke zwischen Rust und Python}
\label{subsec:rust-runner}

Rust und Python sind nicht über eine Sprachbindung wie PyO3 gekoppelt, sondern
lose über einen Subprozess-Aufruf und eine JSON-Datei verbunden. Die Funktion
\texttt{run\_rust\_and\_load\_json} kapselt diese Brücke. Sie prüft zunächst, ob
die Binary existiert, und bestimmt das Arbeitsverzeichnis anhand der
\texttt{input\_config.json}. Eine alte Ausgabedatei wird vorher gelöscht, damit
kein veraltetes Ergebnis eingelesen wird. Nach dem Lauf wertet die Funktion den
Rückgabecode aus; schlägt Rust fehl, landen \texttt{stdout} und \texttt{stderr} in
der Fehlermeldung und erleichtern die Diagnose. Der Vorteil der losen Kopplung
liegt im verteilten Betrieb: Beide Dienste lassen sich getrennt kompilieren,
testen und auf mehrere Cluster-Knoten verteilen.

\subsection{Signalverarbeitung mit SciPy}
\label{subsec:scipy}

Das Modul \texttt{convolution.py} bildet den Kern des Python-Dienstes. Es wandelt
die rohen Verzögerungs- und Druckdaten in eine abtastgenaue Impulsantwort um und
faltet damit das Trockensignal.

\subsubsection{Diskretisierung der Impulsantwort}
\label{subsubsec:build-ir}

Die Rust-Engine liefert eine endliche Menge diskreter Strahl-Treffer, jeweils
bestehend aus einer Ankunftsverzögerung $\tau_j$ (in Sekunden) und einem skalaren
Restdruck $p_j$. Die Funktion \texttt{build\_impulse\_response} quantisiert jeden
Treffer auf den nächstgelegenen Abtastindex $i_j = \operatorname{round}(\tau_j
\cdot f_s)$ und akkumuliert gleichzeitige Treffer zur diskreten Impulsantwort:
\[
    h[i] \;=\; \sum_{j} p_j \cdot \bigl[\,i_j = i\,\bigr].
\]
Die Akkumulation erfolgt mit \texttt{numpy.add.at}. Anders als die einfache
indizierte Zuweisung, die bei doppelten Indizes nur den letzten Wert behält,
summiert \texttt{add.at} gleichzeitige Treffer korrekt auf. Vor der Berechnung
prüft die Funktion die Eingaben an der Systemgrenze: übereinstimmende
Array-Formen, Eindimensionalität, nicht-negative Verzögerungen und eine positive
Abtastrate.

\begin{lstlisting}[style=cemdsp, caption={Diskretisierung der Ray-Hits zur Impulsantwort.}]
indices = np.round(delays * sample_rate).astype(np.int64)
ir_length = int(indices.max()) + 1
ir = np.zeros(ir_length, dtype=np.float64)
np.add.at(ir, indices, pressures)   # akkumuliert gleichzeitige Treffer
return ir
\end{lstlisting}

\subsubsection{Faltung mit \texttt{oaconvolve}}
\label{subsubsec:oaconvolve}

Die eigentliche Faltung nutzt \texttt{scipy.signal.oaconvolve}, die
Overlap-Add-basierte FFT-Faltung aus Abschnitt~2.3. Der Modus \texttt{"full"}
liefert das vollständige Ergebnis der Länge $N + M - 1$ und bewahrt damit den
kompletten Nachhall. Mehrkanaliges (etwa Stereo-)Audio wird kanalweise gefaltet
und danach wieder zusammengesetzt:

\begin{lstlisting}[style=cemdsp, caption={Mono- und Mehrkanal-Faltung.}]
if dry_audio.ndim == 1:
    wet = signal.oaconvolve(dry_audio, ir, mode="full")
elif dry_audio.ndim == 2:
    channels = [
        signal.oaconvolve(dry_audio[:, c], ir, mode="full")
        for c in range(dry_audio.shape[1])
    ]
    wet = np.stack(channels, axis=1)
return wet.astype(np.float32)
\end{lstlisting}

Vor der Faltung prüft \texttt{convolve} das Schema von \texttt{ir\_data} und
vergleicht die Abtastraten von Audio und Impulsantwort. Das Ergebnis gibt die
Funktion roh und unnormalisiert zurück. Die Lautstärke regelt erst der Export
(Abschnitt~\ref{subsec:rendering}); so bleibt die Faltung frei von Pegelfragen.

\subsubsection{Frequenzabhängige Materialdämpfung}
\label{subsubsec:material}

Reale Materialien dämpfen Schall frequenzabhängig. Ein Teppich schluckt hohe
Frequenzen viel stärker als tiefe, Beton reflektiert dagegen fast breitbandig.
Die Rust-Engine reduziert die Absorption aber schon auf einen einzigen skalaren
Restdruck $p_j$ pro Treffer, sodass die spektrale Information an dieser
Systemgrenze verloren geht. Damit der Nachhall trotzdem frequenzabhängig wird,
ergänzt \texttt{convolution.py} ein Absorptionsmodell im Frequenzbereich.

Grundlage sind die Oktavband-Absorptionskoeffizienten $\alpha_b$ aus der Datei
\texttt{materials.json}, die auch die Rust-Engine speist. Physik und
Signalverarbeitung greifen damit auf dieselbe Datenbasis zu (single source of
truth). Jedes Material ist über sechs Oktavbänder mit den Mittenfrequenzen
$125,\,250,\,500,\,1000,\,2000$ und $4000$~Hz beschrieben; ein hoher Koeffizient
bedeutet starke Dämpfung im jeweiligen Band.

Das Modell beruht auf einer einfachen Überlegung: Ein Tap, der zum Zeitpunkt $t$
eintrifft, hat näherungsweise $n(t) = t / \tau_{\mathrm{mfp}}$ Reflexionen hinter
sich, wobei $\tau_{\mathrm{mfp}}$ die mittlere Zeit zwischen zwei Reflexionen
angibt (die mittlere freie Weglänge in Zeiteinheiten). Nach $n$ Reflexionen
verbleibt in jedem Oktavband $b$ der Energieanteil
\[
    g_b(t) \;=\; \bigl(1 - \alpha_b\bigr)^{\,n(t)} \;=\; \bigl(1 - \alpha_b\bigr)^{\,t / \tau_{\mathrm{mfp}}}.
\]
Bänder mit hohem $\alpha_b$, etwa die Höhen bei Teppich, klingen dadurch deutlich
schneller ab als solche mit niedrigem $\alpha_b$. Der Nachhall wird über die Zeit
zunehmend dunkler.

Im Code setzt eine Oktavband-Filterbank dieses Modell um. Die Impulsantwort wird
über Butterworth-Bandpässe (nullphasig via \texttt{sosfiltfilt}) in ihre sechs
Bänder zerlegt, jedes Band mit der zeitabhängigen Hüllkurve $g_b(t)$ gewichtet und
die Bänder anschließend wieder aufsummiert.

\begin{lstlisting}[style=cemdsp, caption={Frequenzabhaengige Material-Daempfung der Impulsantwort.}]
centers, materials = _load_materials(str(_MATERIALS_JSON))
nyquist = sample_rate / 2.0
reflections = (np.arange(ir.size) / sample_rate) / mean_free_path_s

shaped = np.zeros_like(ir)
for (low, high), alpha in zip(_octave_band_edges(centers), materials[material]):
    band = _bandpass(ir, low, high, nyquist)      # ein Oktavband isolieren
    gain = np.power(max(1e-6, 1.0 - alpha), reflections)
    shaped += band * gain                          # zeitabhaengig daempfen
return shaped
\end{lstlisting}

\paragraph{Verifikation.} Ein reproduzierbarer Selbsttest prüft das Modell. Er
dämpft dieselbe dichte Impulsantwort einmal mit \texttt{carpet} und einmal mit
\texttt{concrete} und misst jeweils den Energieanteil oberhalb von $2$~kHz. Bei
Teppich bleiben davon nur $0{,}39$ übrig, bei Beton $0{,}71$; der Teppich schluckt
die Höhen also deutlich stärker. Kippt diese Reihenfolge, schlägt der Test fehl.

\paragraph{Grenzen des Modells.} Das Verfahren modelliert ein globales
Raummaterial, nicht die tatsächliche Materialfolge eines einzelnen Strahlpfads.
Eine Trennung nach Oberfläche, also die Aussage „dieser Tap stammt von Teppich,
jener von Beton", ist im Python-Dienst nicht möglich, solange die Rust-Engine nur
einen skalaren Druck pro Treffer liefert. Dafür müsste die Engine schon bandweise
Drücke ($p_{j,b}$ statt $p_j$) exportieren. Kapitel~7 führt das als mögliche
Erweiterung.

\subsection{Audio-Rendering und Export}
\label{subsec:rendering}

Das Ein- und Auslesen der Audiodateien sowie die Normalisierung gegen Clipping
übernimmt \texttt{wav\_handler.py}; die Details stehen in der Teildokumentation zum
Audio-I/O. Für den DSP-Teil zählt nur der Vertrag zwischen den Stufen.
\texttt{convolve} liefert ein rohes, womöglich übersteuertes Signal;
\texttt{write\_wav} misst dessen maximale Amplitude und normalisiert erst, wenn
sie den Bereich $[-1{,}0,\,1{,}0]$ verlässt. So bleibt Clipping aus, und die
Faltung muss sich nie um Pegel kümmern.

%   \end{document}
% =====================================================================

% Listing-Stil fuer dieses Kapitel (nur listings-Kern, keine Farben noetig)
\lstdefinestyle{cemdsp}{
    language=Python,
    basicstyle=\ttfamily\small,
    showstringspaces=false,
    breaklines=true,
    numbers=left,
    numberstyle=\tiny,
    frame=single,
    tabsize=4
}

%======================================================================
\section{Digitale Signalverarbeitung und Orchestrierung (Python)}
\label{sec:python-dsp}

Die Rust-Engine berechnet die Physik der Schallausbreitung, der Python-Dienst
übernimmt die Signalverarbeitung und die Orchestrierung der Pipeline. Warum die
Physik in Rust und die Signalverarbeitung in Python mit NumPy und SciPy liegt,
ist in Kapitel~3 begründet. Die signaltheoretischen Grundlagen, also
Impulsantwort, Faltungsintegral und schnelle Faltung über die FFT, behandelt
Abschnitt~2.3. Dieses Kapitel zeigt, wie diese Grundlagen im Code umgesetzt sind,
und folgt dabei dem Datenfluss durch die Pipeline.

\subsection{Pipeline-Management durch \texttt{main.py}}
\label{subsec:pipeline}

\texttt{main.py} führt die vier Verarbeitungsschritte nacheinander aus. Tritt in
einem Schritt ein Fehler auf, bricht die Pipeline sofort mit einer Meldung ab,
statt mit ungültigen Zwischendaten weiterzurechnen (fail-fast).

\begin{enumerate}
    \item \textbf{Physik-Engine starten:} \texttt{run\_rust\_and\_load\_json} ruft
    die vorkompilierte Rust-Binary als Subprozess auf und lädt deren Ausgabe
    \texttt{ir\_output.json} als Python-Dictionary (\texttt{ir\_data}).
    \item \textbf{Trockenaudio laden:} \texttt{read\_wav} liest die
    Eingangsdatei und liefert Audiodaten samt Abtastrate.
    \item \textbf{Faltung:} \texttt{convolve} wendet die aus \texttt{ir\_data}
    erzeugte Impulsantwort auf das Trockensignal an.
    \item \textbf{Export:} \texttt{write\_wav} normalisiert das Ergebnis bei
    Bedarf und schreibt die fertige Ausgabedatei.
\end{enumerate}

Der Binary-Name wird plattformabhängig gewählt: unter Windows mit der Endung
\texttt{.exe}, unter macOS und Linux ohne. Stimmt die Abtastrate des Trockenaudios
nicht mit der Rate der Engine überein, gibt \texttt{main.py} eine Warnung aus,
weil der Hall sonst verstimmt klingen kann.

\subsection{Die Prozessbrücke zwischen Rust und Python}
\label{subsec:rust-runner}

Rust und Python sind nicht über eine Sprachbindung wie PyO3 gekoppelt, sondern
lose über einen Subprozess-Aufruf und eine JSON-Datei verbunden. Die Funktion
\texttt{run\_rust\_and\_load\_json} kapselt diese Brücke. Sie prüft zunächst, ob
die Binary existiert, und bestimmt das Arbeitsverzeichnis anhand der
\texttt{input\_config.json}. Eine alte Ausgabedatei wird vorher gelöscht, damit
kein veraltetes Ergebnis eingelesen wird. Nach dem Lauf wertet die Funktion den
Rückgabecode aus; schlägt Rust fehl, landen \texttt{stdout} und \texttt{stderr} in
der Fehlermeldung und erleichtern die Diagnose. Der Vorteil der losen Kopplung
liegt im verteilten Betrieb: Beide Dienste lassen sich getrennt kompilieren,
testen und auf mehrere Cluster-Knoten verteilen.

\subsection{Signalverarbeitung mit SciPy}
\label{subsec:scipy}

Das Modul \texttt{convolution.py} bildet den Kern des Python-Dienstes. Es wandelt
die rohen Verzögerungs- und Druckdaten in eine abtastgenaue Impulsantwort um und
faltet damit das Trockensignal.

\subsubsection{Diskretisierung der Impulsantwort}
\label{subsubsec:build-ir}

Die Rust-Engine liefert eine endliche Menge diskreter Strahl-Treffer, jeweils
bestehend aus einer Ankunftsverzögerung $\tau_j$ (in Sekunden) und einem skalaren
Restdruck $p_j$. Die Funktion \texttt{build\_impulse\_response} quantisiert jeden
Treffer auf den nächstgelegenen Abtastindex $i_j = \operatorname{round}(\tau_j
\cdot f_s)$ und akkumuliert gleichzeitige Treffer zur diskreten Impulsantwort:
\[
    h[i] \;=\; \sum_{j} p_j \cdot \bigl[\,i_j = i\,\bigr].
\]
Die Akkumulation erfolgt mit \texttt{numpy.add.at}. Anders als die einfache
indizierte Zuweisung, die bei doppelten Indizes nur den letzten Wert behält,
summiert \texttt{add.at} gleichzeitige Treffer korrekt auf. Vor der Berechnung
prüft die Funktion die Eingaben an der Systemgrenze: übereinstimmende
Array-Formen, Eindimensionalität, nicht-negative Verzögerungen und eine positive
Abtastrate.

\begin{lstlisting}[style=cemdsp, caption={Diskretisierung der Ray-Hits zur Impulsantwort.}]
indices = np.round(delays * sample_rate).astype(np.int64)
ir_length = int(indices.max()) + 1
ir = np.zeros(ir_length, dtype=np.float64)
np.add.at(ir, indices, pressures)   # akkumuliert gleichzeitige Treffer
return ir
\end{lstlisting}

\subsubsection{Faltung mit \texttt{oaconvolve}}
\label{subsubsec:oaconvolve}

Die eigentliche Faltung nutzt \texttt{scipy.signal.oaconvolve}, die
Overlap-Add-basierte FFT-Faltung aus Abschnitt~2.3. Der Modus \texttt{"full"}
liefert das vollständige Ergebnis der Länge $N + M - 1$ und bewahrt damit den
kompletten Nachhall. Mehrkanaliges (etwa Stereo-)Audio wird kanalweise gefaltet
und danach wieder zusammengesetzt:

\begin{lstlisting}[style=cemdsp, caption={Mono- und Mehrkanal-Faltung.}]
if dry_audio.ndim == 1:
    wet = signal.oaconvolve(dry_audio, ir, mode="full")
elif dry_audio.ndim == 2:
    channels = [
        signal.oaconvolve(dry_audio[:, c], ir, mode="full")
        for c in range(dry_audio.shape[1])
    ]
    wet = np.stack(channels, axis=1)
return wet.astype(np.float32)
\end{lstlisting}

Vor der Faltung prüft \texttt{convolve} das Schema von \texttt{ir\_data} und
vergleicht die Abtastraten von Audio und Impulsantwort. Das Ergebnis gibt die
Funktion roh und unnormalisiert zurück. Die Lautstärke regelt erst der Export
(Abschnitt~\ref{subsec:rendering}); so bleibt die Faltung frei von Pegelfragen.

\subsubsection{Frequenzabhängige Materialdämpfung}
\label{subsubsec:material}

Reale Materialien dämpfen Schall frequenzabhängig. Ein Teppich schluckt hohe
Frequenzen viel stärker als tiefe, Beton reflektiert dagegen fast breitbandig.
Die Rust-Engine reduziert die Absorption aber schon auf einen einzigen skalaren
Restdruck $p_j$ pro Treffer, sodass die spektrale Information an dieser
Systemgrenze verloren geht. Damit der Nachhall trotzdem frequenzabhängig wird,
ergänzt \texttt{convolution.py} ein Absorptionsmodell im Frequenzbereich.

Grundlage sind die Oktavband-Absorptionskoeffizienten $\alpha_b$ aus der Datei
\texttt{materials.json}, die auch die Rust-Engine speist. Physik und
Signalverarbeitung greifen damit auf dieselbe Datenbasis zu (single source of
truth). Jedes Material ist über sechs Oktavbänder mit den Mittenfrequenzen
$125,\,250,\,500,\,1000,\,2000$ und $4000$~Hz beschrieben; ein hoher Koeffizient
bedeutet starke Dämpfung im jeweiligen Band.

Das Modell beruht auf einer einfachen Überlegung: Ein Tap, der zum Zeitpunkt $t$
eintrifft, hat näherungsweise $n(t) = t / \tau_{\mathrm{mfp}}$ Reflexionen hinter
sich, wobei $\tau_{\mathrm{mfp}}$ die mittlere Zeit zwischen zwei Reflexionen
angibt (die mittlere freie Weglänge in Zeiteinheiten). Nach $n$ Reflexionen
verbleibt in jedem Oktavband $b$ der Energieanteil
\[
    g_b(t) \;=\; \bigl(1 - \alpha_b\bigr)^{\,n(t)} \;=\; \bigl(1 - \alpha_b\bigr)^{\,t / \tau_{\mathrm{mfp}}}.
\]
Bänder mit hohem $\alpha_b$, etwa die Höhen bei Teppich, klingen dadurch deutlich
schneller ab als solche mit niedrigem $\alpha_b$. Der Nachhall wird über die Zeit
zunehmend dunkler.

Im Code setzt eine Oktavband-Filterbank dieses Modell um. Die Impulsantwort wird
über Butterworth-Bandpässe (nullphasig via \texttt{sosfiltfilt}) in ihre sechs
Bänder zerlegt, jedes Band mit der zeitabhängigen Hüllkurve $g_b(t)$ gewichtet und
die Bänder anschließend wieder aufsummiert.

\begin{lstlisting}[style=cemdsp, caption={Frequenzabhaengige Material-Daempfung der Impulsantwort.}]
centers, materials = _load_materials(str(_MATERIALS_JSON))
nyquist = sample_rate / 2.0
reflections = (np.arange(ir.size) / sample_rate) / mean_free_path_s

shaped = np.zeros_like(ir)
for (low, high), alpha in zip(_octave_band_edges(centers), materials[material]):
    band = _bandpass(ir, low, high, nyquist)      # ein Oktavband isolieren
    gain = np.power(max(1e-6, 1.0 - alpha), reflections)
    shaped += band * gain                          # zeitabhaengig daempfen
return shaped
\end{lstlisting}

\paragraph{Verifikation.} Ein reproduzierbarer Selbsttest prüft das Modell. Er
dämpft dieselbe dichte Impulsantwort einmal mit \texttt{carpet} und einmal mit
\texttt{concrete} und misst jeweils den Energieanteil oberhalb von $2$~kHz. Bei
Teppich bleiben davon nur $0{,}39$ übrig, bei Beton $0{,}71$; der Teppich schluckt
die Höhen also deutlich stärker. Kippt diese Reihenfolge, schlägt der Test fehl.

\paragraph{Grenzen des Modells.} Das Verfahren modelliert ein globales
Raummaterial, nicht die tatsächliche Materialfolge eines einzelnen Strahlpfads.
Eine Trennung nach Oberfläche, also die Aussage „dieser Tap stammt von Teppich,
jener von Beton", ist im Python-Dienst nicht möglich, solange die Rust-Engine nur
einen skalaren Druck pro Treffer liefert. Dafür müsste die Engine schon bandweise
Drücke ($p_{j,b}$ statt $p_j$) exportieren. Kapitel~7 führt das als mögliche
Erweiterung.

\subsection{Audio-Rendering und Export}
\label{subsec:rendering}

Das Ein- und Auslesen der Audiodateien sowie die Normalisierung gegen Clipping
übernimmt \texttt{wav\_handler.py}; die Details stehen in der Teildokumentation zum
Audio-I/O. Für den DSP-Teil zählt nur der Vertrag zwischen den Stufen.
\texttt{convolve} liefert ein rohes, womöglich übersteuertes Signal;
\texttt{write\_wav} misst dessen maximale Amplitude und normalisiert erst, wenn
sie den Bereich $[-1{,}0,\,1{,}0]$ verlässt. So bleibt Clipping aus, und die
Faltung muss sich nie um Pegel kümmern.

%   \end{document}
% =====================================================================

% Listing-Stil fuer dieses Kapitel (nur listings-Kern, keine Farben noetig)
\lstdefinestyle{cemdsp}{
    language=Python,
    basicstyle=\ttfamily\small,
    showstringspaces=false,
    breaklines=true,
    numbers=left,
    numberstyle=\tiny,
    frame=single,
    tabsize=4
}

%======================================================================
\section{Digitale Signalverarbeitung und Orchestrierung (Python)}
\label{sec:python-dsp}

Die Rust-Engine berechnet die Physik der Schallausbreitung, der Python-Dienst
übernimmt die Signalverarbeitung und die Orchestrierung der Pipeline. Warum die
Physik in Rust und die Signalverarbeitung in Python mit NumPy und SciPy liegt,
ist in Kapitel~3 begründet. Die signaltheoretischen Grundlagen, also
Impulsantwort, Faltungsintegral und schnelle Faltung über die FFT, behandelt
Abschnitt~2.3. Dieses Kapitel zeigt, wie diese Grundlagen im Code umgesetzt sind,
und folgt dabei dem Datenfluss durch die Pipeline.

\subsection{Pipeline-Management durch \texttt{main.py}}
\label{subsec:pipeline}

\texttt{main.py} führt die vier Verarbeitungsschritte nacheinander aus. Tritt in
einem Schritt ein Fehler auf, bricht die Pipeline sofort mit einer Meldung ab,
statt mit ungültigen Zwischendaten weiterzurechnen (fail-fast).

\begin{enumerate}
    \item \textbf{Physik-Engine starten:} \texttt{run\_rust\_and\_load\_json} ruft
    die vorkompilierte Rust-Binary als Subprozess auf und lädt deren Ausgabe
    \texttt{ir\_output.json} als Python-Dictionary (\texttt{ir\_data}).
    \item \textbf{Trockenaudio laden:} \texttt{read\_wav} liest die
    Eingangsdatei und liefert Audiodaten samt Abtastrate.
    \item \textbf{Faltung:} \texttt{convolve} wendet die aus \texttt{ir\_data}
    erzeugte Impulsantwort auf das Trockensignal an.
    \item \textbf{Export:} \texttt{write\_wav} normalisiert das Ergebnis bei
    Bedarf und schreibt die fertige Ausgabedatei.
\end{enumerate}

Der Binary-Name wird plattformabhängig gewählt: unter Windows mit der Endung
\texttt{.exe}, unter macOS und Linux ohne. Stimmt die Abtastrate des Trockenaudios
nicht mit der Rate der Engine überein, gibt \texttt{main.py} eine Warnung aus,
weil der Hall sonst verstimmt klingen kann.

\subsection{Die Prozessbrücke zwischen Rust und Python}
\label{subsec:rust-runner}

Rust und Python sind nicht über eine Sprachbindung wie PyO3 gekoppelt, sondern
lose über einen Subprozess-Aufruf und eine JSON-Datei verbunden. Die Funktion
\texttt{run\_rust\_and\_load\_json} kapselt diese Brücke. Sie prüft zunächst, ob
die Binary existiert, und bestimmt das Arbeitsverzeichnis anhand der
\texttt{input\_config.json}. Eine alte Ausgabedatei wird vorher gelöscht, damit
kein veraltetes Ergebnis eingelesen wird. Nach dem Lauf wertet die Funktion den
Rückgabecode aus; schlägt Rust fehl, landen \texttt{stdout} und \texttt{stderr} in
der Fehlermeldung und erleichtern die Diagnose. Der Vorteil der losen Kopplung
liegt im verteilten Betrieb: Beide Dienste lassen sich getrennt kompilieren,
testen und auf mehrere Cluster-Knoten verteilen.

\subsection{Signalverarbeitung mit SciPy}
\label{subsec:scipy}

Das Modul \texttt{convolution.py} bildet den Kern des Python-Dienstes. Es wandelt
die rohen Verzögerungs- und Druckdaten in eine abtastgenaue Impulsantwort um und
faltet damit das Trockensignal.

\subsubsection{Diskretisierung der Impulsantwort}
\label{subsubsec:build-ir}

Die Rust-Engine liefert eine endliche Menge diskreter Strahl-Treffer, jeweils
bestehend aus einer Ankunftsverzögerung $\tau_j$ (in Sekunden) und einem skalaren
Restdruck $p_j$. Die Funktion \texttt{build\_impulse\_response} quantisiert jeden
Treffer auf den nächstgelegenen Abtastindex $i_j = \operatorname{round}(\tau_j
\cdot f_s)$ und akkumuliert gleichzeitige Treffer zur diskreten Impulsantwort:
\[
    h[i] \;=\; \sum_{j} p_j \cdot \bigl[\,i_j = i\,\bigr].
\]
Die Akkumulation erfolgt mit \texttt{numpy.add.at}. Anders als die einfache
indizierte Zuweisung, die bei doppelten Indizes nur den letzten Wert behält,
summiert \texttt{add.at} gleichzeitige Treffer korrekt auf. Vor der Berechnung
prüft die Funktion die Eingaben an der Systemgrenze: übereinstimmende
Array-Formen, Eindimensionalität, nicht-negative Verzögerungen und eine positive
Abtastrate.

\begin{lstlisting}[style=cemdsp, caption={Diskretisierung der Ray-Hits zur Impulsantwort.}]
indices = np.round(delays * sample_rate).astype(np.int64)
ir_length = int(indices.max()) + 1
ir = np.zeros(ir_length, dtype=np.float64)
np.add.at(ir, indices, pressures)   # akkumuliert gleichzeitige Treffer
return ir
\end{lstlisting}

\subsubsection{Faltung mit \texttt{oaconvolve}}
\label{subsubsec:oaconvolve}

Die eigentliche Faltung nutzt \texttt{scipy.signal.oaconvolve}, die
Overlap-Add-basierte FFT-Faltung aus Abschnitt~2.3. Der Modus \texttt{"full"}
liefert das vollständige Ergebnis der Länge $N + M - 1$ und bewahrt damit den
kompletten Nachhall. Mehrkanaliges (etwa Stereo-)Audio wird kanalweise gefaltet
und danach wieder zusammengesetzt:

\begin{lstlisting}[style=cemdsp, caption={Mono- und Mehrkanal-Faltung.}]
if dry_audio.ndim == 1:
    wet = signal.oaconvolve(dry_audio, ir, mode="full")
elif dry_audio.ndim == 2:
    channels = [
        signal.oaconvolve(dry_audio[:, c], ir, mode="full")
        for c in range(dry_audio.shape[1])
    ]
    wet = np.stack(channels, axis=1)
return wet.astype(np.float32)
\end{lstlisting}

Vor der Faltung prüft \texttt{convolve} das Schema von \texttt{ir\_data} und
vergleicht die Abtastraten von Audio und Impulsantwort. Das Ergebnis gibt die
Funktion roh und unnormalisiert zurück. Die Lautstärke regelt erst der Export
(Abschnitt~\ref{subsec:rendering}); so bleibt die Faltung frei von Pegelfragen.

\subsubsection{Frequenzabhängige Materialdämpfung}
\label{subsubsec:material}

Reale Materialien dämpfen Schall frequenzabhängig. Ein Teppich schluckt hohe
Frequenzen viel stärker als tiefe, Beton reflektiert dagegen fast breitbandig.
Die Rust-Engine reduziert die Absorption aber schon auf einen einzigen skalaren
Restdruck $p_j$ pro Treffer, sodass die spektrale Information an dieser
Systemgrenze verloren geht. Damit der Nachhall trotzdem frequenzabhängig wird,
ergänzt \texttt{convolution.py} ein Absorptionsmodell im Frequenzbereich.

Grundlage sind die Oktavband-Absorptionskoeffizienten $\alpha_b$ aus der Datei
\texttt{materials.json}, die auch die Rust-Engine speist. Physik und
Signalverarbeitung greifen damit auf dieselbe Datenbasis zu (single source of
truth). Jedes Material ist über sechs Oktavbänder mit den Mittenfrequenzen
$125,\,250,\,500,\,1000,\,2000$ und $4000$~Hz beschrieben; ein hoher Koeffizient
bedeutet starke Dämpfung im jeweiligen Band.

Das Modell beruht auf einer einfachen Überlegung: Ein Tap, der zum Zeitpunkt $t$
eintrifft, hat näherungsweise $n(t) = t / \tau_{\mathrm{mfp}}$ Reflexionen hinter
sich, wobei $\tau_{\mathrm{mfp}}$ die mittlere Zeit zwischen zwei Reflexionen
angibt (die mittlere freie Weglänge in Zeiteinheiten). Nach $n$ Reflexionen
verbleibt in jedem Oktavband $b$ der Energieanteil
\[
    g_b(t) \;=\; \bigl(1 - \alpha_b\bigr)^{\,n(t)} \;=\; \bigl(1 - \alpha_b\bigr)^{\,t / \tau_{\mathrm{mfp}}}.
\]
Bänder mit hohem $\alpha_b$, etwa die Höhen bei Teppich, klingen dadurch deutlich
schneller ab als solche mit niedrigem $\alpha_b$. Der Nachhall wird über die Zeit
zunehmend dunkler.

Im Code setzt eine Oktavband-Filterbank dieses Modell um. Die Impulsantwort wird
über Butterworth-Bandpässe (nullphasig via \texttt{sosfiltfilt}) in ihre sechs
Bänder zerlegt, jedes Band mit der zeitabhängigen Hüllkurve $g_b(t)$ gewichtet und
die Bänder anschließend wieder aufsummiert.

\begin{lstlisting}[style=cemdsp, caption={Frequenzabhaengige Material-Daempfung der Impulsantwort.}]
centers, materials = _load_materials(str(_MATERIALS_JSON))
nyquist = sample_rate / 2.0
reflections = (np.arange(ir.size) / sample_rate) / mean_free_path_s

shaped = np.zeros_like(ir)
for (low, high), alpha in zip(_octave_band_edges(centers), materials[material]):
    band = _bandpass(ir, low, high, nyquist)      # ein Oktavband isolieren
    gain = np.power(max(1e-6, 1.0 - alpha), reflections)
    shaped += band * gain                          # zeitabhaengig daempfen
return shaped
\end{lstlisting}

\paragraph{Verifikation.} Ein reproduzierbarer Selbsttest prüft das Modell. Er
dämpft dieselbe dichte Impulsantwort einmal mit \texttt{carpet} und einmal mit
\texttt{concrete} und misst jeweils den Energieanteil oberhalb von $2$~kHz. Bei
Teppich bleiben davon nur $0{,}39$ übrig, bei Beton $0{,}71$; der Teppich schluckt
die Höhen also deutlich stärker. Kippt diese Reihenfolge, schlägt der Test fehl.

\paragraph{Grenzen des Modells.} Das Verfahren modelliert ein globales
Raummaterial, nicht die tatsächliche Materialfolge eines einzelnen Strahlpfads.
Eine Trennung nach Oberfläche, also die Aussage „dieser Tap stammt von Teppich,
jener von Beton", ist im Python-Dienst nicht möglich, solange die Rust-Engine nur
einen skalaren Druck pro Treffer liefert. Dafür müsste die Engine schon bandweise
Drücke ($p_{j,b}$ statt $p_j$) exportieren. Kapitel~7 führt das als mögliche
Erweiterung.

\subsection{Audio-Rendering und Export}
\label{subsec:rendering}

Das Ein- und Auslesen der Audiodateien sowie die Normalisierung gegen Clipping
übernimmt \texttt{wav\_handler.py}; die Details stehen in der Teildokumentation zum
Audio-I/O. Für den DSP-Teil zählt nur der Vertrag zwischen den Stufen.
\texttt{convolve} liefert ein rohes, womöglich übersteuertes Signal;
\texttt{write\_wav} misst dessen maximale Amplitude und normalisiert erst, wenn
sie den Bereich $[-1{,}0,\,1{,}0]$ verlässt. So bleibt Clipping aus, und die
Faltung muss sich nie um Pegel kümmern.

%   \end{document}
% =====================================================================

% Listing-Stil fuer dieses Kapitel (nur listings-Kern, keine Farben noetig)
\lstdefinestyle{cemdsp}{
    language=Python,
    basicstyle=\ttfamily\small,
    showstringspaces=false,
    breaklines=true,
    numbers=left,
    numberstyle=\tiny,
    frame=single,
    tabsize=4
}

%======================================================================
\section{Digitale Signalverarbeitung und Orchestrierung (Python)}
\label{sec:python-dsp}

Die Rust-Engine berechnet die Physik der Schallausbreitung, der Python-Dienst
übernimmt die Signalverarbeitung und die Orchestrierung der Pipeline. Warum die
Physik in Rust und die Signalverarbeitung in Python mit NumPy und SciPy liegt,
ist in Kapitel~3 begründet. Die signaltheoretischen Grundlagen, also
Impulsantwort, Faltungsintegral und schnelle Faltung über die FFT, behandelt
Abschnitt~2.3. Dieses Kapitel zeigt, wie diese Grundlagen im Code umgesetzt sind,
und folgt dabei dem Datenfluss durch die Pipeline.

\subsection{Pipeline-Management durch \texttt{main.py}}
\label{subsec:pipeline}

\texttt{main.py} führt die vier Verarbeitungsschritte nacheinander aus. Tritt in
einem Schritt ein Fehler auf, bricht die Pipeline sofort mit einer Meldung ab,
statt mit ungültigen Zwischendaten weiterzurechnen (fail-fast).

\begin{enumerate}
    \item \textbf{Physik-Engine starten:} \texttt{run\_rust\_and\_load\_json} ruft
    die vorkompilierte Rust-Binary als Subprozess auf und lädt deren Ausgabe
    \texttt{ir\_output.json} als Python-Dictionary (\texttt{ir\_data}).
    \item \textbf{Trockenaudio laden:} \texttt{read\_wav} liest die
    Eingangsdatei und liefert Audiodaten samt Abtastrate.
    \item \textbf{Faltung:} \texttt{convolve} wendet die aus \texttt{ir\_data}
    erzeugte Impulsantwort auf das Trockensignal an.
    \item \textbf{Export:} \texttt{write\_wav} normalisiert das Ergebnis bei
    Bedarf und schreibt die fertige Ausgabedatei.
\end{enumerate}

Der Binary-Name wird plattformabhängig gewählt: unter Windows mit der Endung
\texttt{.exe}, unter macOS und Linux ohne. Stimmt die Abtastrate des Trockenaudios
nicht mit der Rate der Engine überein, gibt \texttt{main.py} eine Warnung aus,
weil der Hall sonst verstimmt klingen kann.

\subsection{Die Prozessbrücke zwischen Rust und Python}
\label{subsec:rust-runner}

Rust und Python sind nicht über eine Sprachbindung wie PyO3 gekoppelt, sondern
lose über einen Subprozess-Aufruf und eine JSON-Datei verbunden. Die Funktion
\texttt{run\_rust\_and\_load\_json} kapselt diese Brücke. Sie prüft zunächst, ob
die Binary existiert, und bestimmt das Arbeitsverzeichnis anhand der
\texttt{input\_config.json}. Eine alte Ausgabedatei wird vorher gelöscht, damit
kein veraltetes Ergebnis eingelesen wird. Nach dem Lauf wertet die Funktion den
Rückgabecode aus; schlägt Rust fehl, landen \texttt{stdout} und \texttt{stderr} in
der Fehlermeldung und erleichtern die Diagnose. Der Vorteil der losen Kopplung
liegt im verteilten Betrieb: Beide Dienste lassen sich getrennt kompilieren,
testen und auf mehrere Cluster-Knoten verteilen.

\subsection{Signalverarbeitung mit SciPy}
\label{subsec:scipy}

Das Modul \texttt{convolution.py} bildet den Kern des Python-Dienstes. Es wandelt
die rohen Verzögerungs- und Druckdaten in eine abtastgenaue Impulsantwort um und
faltet damit das Trockensignal.

\subsubsection{Diskretisierung der Impulsantwort}
\label{subsubsec:build-ir}

Die Rust-Engine liefert eine endliche Menge diskreter Strahl-Treffer, jeweils
bestehend aus einer Ankunftsverzögerung $\tau_j$ (in Sekunden) und einem skalaren
Restdruck $p_j$. Die Funktion \texttt{build\_impulse\_response} quantisiert jeden
Treffer auf den nächstgelegenen Abtastindex $i_j = \operatorname{round}(\tau_j
\cdot f_s)$ und akkumuliert gleichzeitige Treffer zur diskreten Impulsantwort:
\[
    h[i] \;=\; \sum_{j} p_j \cdot \bigl[\,i_j = i\,\bigr].
\]
Die Akkumulation erfolgt mit \texttt{numpy.add.at}. Anders als die einfache
indizierte Zuweisung, die bei doppelten Indizes nur den letzten Wert behält,
summiert \texttt{add.at} gleichzeitige Treffer korrekt auf. Vor der Berechnung
prüft die Funktion die Eingaben an der Systemgrenze: übereinstimmende
Array-Formen, Eindimensionalität, nicht-negative Verzögerungen und eine positive
Abtastrate.

\begin{lstlisting}[style=cemdsp, caption={Diskretisierung der Ray-Hits zur Impulsantwort.}]
indices = np.round(delays * sample_rate).astype(np.int64)
ir_length = int(indices.max()) + 1
ir = np.zeros(ir_length, dtype=np.float64)
np.add.at(ir, indices, pressures)   # akkumuliert gleichzeitige Treffer
return ir
\end{lstlisting}

\subsubsection{Faltung mit \texttt{oaconvolve}}
\label{subsubsec:oaconvolve}

Die eigentliche Faltung nutzt \texttt{scipy.signal.oaconvolve}, die
Overlap-Add-basierte FFT-Faltung aus Abschnitt~2.3. Der Modus \texttt{"full"}
liefert das vollständige Ergebnis der Länge $N + M - 1$ und bewahrt damit den
kompletten Nachhall. Mehrkanaliges (etwa Stereo-)Audio wird kanalweise gefaltet
und danach wieder zusammengesetzt:

\begin{lstlisting}[style=cemdsp, caption={Mono- und Mehrkanal-Faltung.}]
if dry_audio.ndim == 1:
    wet = signal.oaconvolve(dry_audio, ir, mode="full")
elif dry_audio.ndim == 2:
    channels = [
        signal.oaconvolve(dry_audio[:, c], ir, mode="full")
        for c in range(dry_audio.shape[1])
    ]
    wet = np.stack(channels, axis=1)
return wet.astype(np.float32)
\end{lstlisting}

Vor der Faltung prüft \texttt{convolve} das Schema von \texttt{ir\_data} und
vergleicht die Abtastraten von Audio und Impulsantwort. Das Ergebnis gibt die
Funktion roh und unnormalisiert zurück. Die Lautstärke regelt erst der Export
(Abschnitt~\ref{subsec:rendering}); so bleibt die Faltung frei von Pegelfragen.

\subsubsection{Frequenzabhängige Materialdämpfung}
\label{subsubsec:material}

Reale Materialien dämpfen Schall frequenzabhängig. Ein Teppich schluckt hohe
Frequenzen viel stärker als tiefe, Beton reflektiert dagegen fast breitbandig.
Die Rust-Engine reduziert die Absorption aber schon auf einen einzigen skalaren
Restdruck $p_j$ pro Treffer, sodass die spektrale Information an dieser
Systemgrenze verloren geht. Damit der Nachhall trotzdem frequenzabhängig wird,
ergänzt \texttt{convolution.py} ein Absorptionsmodell im Frequenzbereich.

Grundlage sind die Oktavband-Absorptionskoeffizienten $\alpha_b$ aus der Datei
\texttt{materials.json}, die auch die Rust-Engine speist. Physik und
Signalverarbeitung greifen damit auf dieselbe Datenbasis zu (single source of
truth). Jedes Material ist über sechs Oktavbänder mit den Mittenfrequenzen
$125,\,250,\,500,\,1000,\,2000$ und $4000$~Hz beschrieben; ein hoher Koeffizient
bedeutet starke Dämpfung im jeweiligen Band.

Das Modell beruht auf einer einfachen Überlegung: Ein Tap, der zum Zeitpunkt $t$
eintrifft, hat näherungsweise $n(t) = t / \tau_{\mathrm{mfp}}$ Reflexionen hinter
sich, wobei $\tau_{\mathrm{mfp}}$ die mittlere Zeit zwischen zwei Reflexionen
angibt (die mittlere freie Weglänge in Zeiteinheiten). Nach $n$ Reflexionen
verbleibt in jedem Oktavband $b$ der Energieanteil
\[
    g_b(t) \;=\; \bigl(1 - \alpha_b\bigr)^{\,n(t)} \;=\; \bigl(1 - \alpha_b\bigr)^{\,t / \tau_{\mathrm{mfp}}}.
\]
Bänder mit hohem $\alpha_b$, etwa die Höhen bei Teppich, klingen dadurch deutlich
schneller ab als solche mit niedrigem $\alpha_b$. Der Nachhall wird über die Zeit
zunehmend dunkler.

Im Code setzt eine Oktavband-Filterbank dieses Modell um. Die Impulsantwort wird
über Butterworth-Bandpässe (nullphasig via \texttt{sosfiltfilt}) in ihre sechs
Bänder zerlegt, jedes Band mit der zeitabhängigen Hüllkurve $g_b(t)$ gewichtet und
die Bänder anschließend wieder aufsummiert.

\begin{lstlisting}[style=cemdsp, caption={Frequenzabhaengige Material-Daempfung der Impulsantwort.}]
centers, materials = _load_materials(str(_MATERIALS_JSON))
nyquist = sample_rate / 2.0
reflections = (np.arange(ir.size) / sample_rate) / mean_free_path_s

shaped = np.zeros_like(ir)
for (low, high), alpha in zip(_octave_band_edges(centers), materials[material]):
    band = _bandpass(ir, low, high, nyquist)      # ein Oktavband isolieren
    gain = np.power(max(1e-6, 1.0 - alpha), reflections)
    shaped += band * gain                          # zeitabhaengig daempfen
return shaped
\end{lstlisting}

\paragraph{Verifikation.} Ein reproduzierbarer Selbsttest prüft das Modell. Er
dämpft dieselbe dichte Impulsantwort einmal mit \texttt{carpet} und einmal mit
\texttt{concrete} und misst jeweils den Energieanteil oberhalb von $2$~kHz. Bei
Teppich bleiben davon nur $0{,}39$ übrig, bei Beton $0{,}71$; der Teppich schluckt
die Höhen also deutlich stärker. Kippt diese Reihenfolge, schlägt der Test fehl.

\paragraph{Grenzen des Modells.} Das Verfahren modelliert ein globales
Raummaterial, nicht die tatsächliche Materialfolge eines einzelnen Strahlpfads.
Eine Trennung nach Oberfläche, also die Aussage „dieser Tap stammt von Teppich,
jener von Beton", ist im Python-Dienst nicht möglich, solange die Rust-Engine nur
einen skalaren Druck pro Treffer liefert. Dafür müsste die Engine schon bandweise
Drücke ($p_{j,b}$ statt $p_j$) exportieren. Kapitel~7 führt das als mögliche
Erweiterung.

\subsection{Audio-Rendering und Export}
\label{subsec:rendering}

Das Ein- und Auslesen der Audiodateien sowie die Normalisierung gegen Clipping
übernimmt \texttt{wav\_handler.py}; die Details stehen in der Teildokumentation zum
Audio-I/O. Für den DSP-Teil zählt nur der Vertrag zwischen den Stufen.
\texttt{convolve} liefert ein rohes, womöglich übersteuertes Signal;
\texttt{write\_wav} misst dessen maximale Amplitude und normalisiert erst, wenn
sie den Bereich $[-1{,}0,\,1{,}0]$ verlässt. So bleibt Clipping aus, und die
Faltung muss sich nie um Pegel kümmern.
